% Options for packages loaded elsewhere
\PassOptionsToPackage{unicode}{hyperref}
\PassOptionsToPackage{hyphens}{url}
%
\documentclass[
]{article}
\usepackage{amsmath,amssymb}
\usepackage{iftex}
\ifPDFTeX
  \usepackage[T1]{fontenc}
  \usepackage[utf8]{inputenc}
  \usepackage{textcomp} % provide euro and other symbols
\else % if luatex or xetex
  \usepackage{unicode-math} % this also loads fontspec
  \defaultfontfeatures{Scale=MatchLowercase}
  \defaultfontfeatures[\rmfamily]{Ligatures=TeX,Scale=1}
\fi
\usepackage{lmodern}
\ifPDFTeX\else
  % xetex/luatex font selection
\fi
% Use upquote if available, for straight quotes in verbatim environments
\IfFileExists{upquote.sty}{\usepackage{upquote}}{}
\IfFileExists{microtype.sty}{% use microtype if available
  \usepackage[]{microtype}
  \UseMicrotypeSet[protrusion]{basicmath} % disable protrusion for tt fonts
}{}
\makeatletter
\@ifundefined{KOMAClassName}{% if non-KOMA class
  \IfFileExists{parskip.sty}{%
    \usepackage{parskip}
  }{% else
    \setlength{\parindent}{0pt}
    \setlength{\parskip}{6pt plus 2pt minus 1pt}}
}{% if KOMA class
  \KOMAoptions{parskip=half}}
\makeatother
\usepackage{xcolor}
\usepackage{longtable,booktabs,array}
\usepackage{calc} % for calculating minipage widths
% Correct order of tables after \paragraph or \subparagraph
\usepackage{etoolbox}
\makeatletter
\patchcmd\longtable{\par}{\if@noskipsec\mbox{}\fi\par}{}{}
\makeatother
% Allow footnotes in longtable head/foot
\IfFileExists{footnotehyper.sty}{\usepackage{footnotehyper}}{\usepackage{footnote}}
\makesavenoteenv{longtable}
\usepackage{graphicx}
\makeatletter
\def\maxwidth{\ifdim\Gin@nat@width>\linewidth\linewidth\else\Gin@nat@width\fi}
\def\maxheight{\ifdim\Gin@nat@height>\textheight\textheight\else\Gin@nat@height\fi}
\makeatother
% Scale images if necessary, so that they will not overflow the page
% margins by default, and it is still possible to overwrite the defaults
% using explicit options in \includegraphics[width, height, ...]{}
\setkeys{Gin}{width=\maxwidth,height=\maxheight,keepaspectratio}
% Set default figure placement to htbp
\makeatletter
\def\fps@figure{htbp}
\makeatother
\setlength{\emergencystretch}{3em} % prevent overfull lines
\providecommand{\tightlist}{%
  \setlength{\itemsep}{0pt}\setlength{\parskip}{0pt}}
\setcounter{secnumdepth}{-\maxdimen} % remove section numbering
% Soporte para caracteres Unicode (modo una columna)
\usepackage{newunicodechar}
\usepackage{array}
\usepackage{longtable}
\usepackage{booktabs}
\usepackage{tabularx}
\usepackage{ragged2e}

% Definir caracteres Unicode problemáticos
\newunicodechar{₀}{\ensuremath{_0}}
\newunicodechar{₁}{\ensuremath{_1}}
\newunicodechar{₂}{\ensuremath{_2}}
\newunicodechar{₃}{\ensuremath{_3}}
\newunicodechar{₄}{\ensuremath{_4}}
\newunicodechar{₅}{\ensuremath{_5}}
\newunicodechar{₆}{\ensuremath{_6}}
\newunicodechar{₇}{\ensuremath{_7}}
\newunicodechar{₈}{\ensuremath{_8}}
\newunicodechar{₉}{\ensuremath{_9}}
\newunicodechar{⁺}{\ensuremath{^+}}
\newunicodechar{⁻}{\ensuremath{^-}}
\newunicodechar{°}{\textdegree}
\newunicodechar{×}{\texttimes}
\newunicodechar{·}{\textperiodcentered}
\newunicodechar{–}{\textendash}
\newunicodechar{—}{\textemdash}
\newunicodechar{λ}{\ensuremath{\lambda}}
\newunicodechar{σ}{\ensuremath{\sigma}}
\newunicodechar{ε}{\ensuremath{\varepsilon}}
\newunicodechar{μ}{\ensuremath{\mu}}
\newunicodechar{χ}{\ensuremath{\chi}}
\newunicodechar{α}{\ensuremath{\alpha}}
\newunicodechar{β}{\ensuremath{\beta}}
\newunicodechar{γ}{\ensuremath{\gamma}}
\newunicodechar{δ}{\ensuremath{\delta}}
\newunicodechar{θ}{\ensuremath{\theta}}
\newunicodechar{φ}{\ensuremath{\phi}}
\newunicodechar{ω}{\ensuremath{\omega}}
\newunicodechar{π}{\ensuremath{\pi}}
\newunicodechar{Σ}{\ensuremath{\Sigma}}
\newunicodechar{Δ}{\ensuremath{\Delta}}
\newunicodechar{∞}{\ensuremath{\infty}}
\newunicodechar{≈}{\ensuremath{\approx}}
\newunicodechar{≠}{\ensuremath{\neq}}
\newunicodechar{≤}{\ensuremath{\leq}}
\newunicodechar{≥}{\ensuremath{\geq}}
\newunicodechar{→}{\ensuremath{\rightarrow}}
\newunicodechar{←}{\ensuremath{\leftarrow}}
\newunicodechar{↔}{\ensuremath{\leftrightarrow}}
\newunicodechar{↑}{\ensuremath{\uparrow}}
\newunicodechar{↓}{\ensuremath{\downarrow}}
\newunicodechar{±}{\ensuremath{\pm}}
\newunicodechar{÷}{\ensuremath{\div}}
\newunicodechar{≪}{\ensuremath{\ll}}
\newunicodechar{≫}{\ensuremath{\gg}}
\newunicodechar{∇}{\ensuremath{\nabla}}
\newunicodechar{∂}{\ensuremath{\partial}}

% Configuración para tablas
\setlength{\LTpre}{\baselineskip}
\setlength{\LTpost}{\baselineskip}
\renewcommand{\arraystretch}{1.2}

% Permitir que las tablas se ajusten al ancho de página
\makeatletter
\def\tabularxcolumn#1{m{#1}}
\makeatother
\ifLuaTeX
  \usepackage{selnolig}  % disable illegal ligatures
\fi
\usepackage{bookmark}
\IfFileExists{xurl.sty}{\usepackage{xurl}}{} % add URL line breaks if available
\urlstyle{same}
\hypersetup{
  hidelinks,
  pdfcreator={LaTeX via pandoc}}

\author{}
\date{}

\begin{document}

\section{En los orígenes de la electrocultura: Un modelado retrodictivo
del electrovegetómetro de Bertholon del siglo XVIII en el régimen
pre-corona}\label{en-los-oruxedgenes-de-la-electrocultura-un-modelado-retrodictivo-del-electrovegetuxf3metro-de-bertholon-del-siglo-xviii-en-el-ruxe9gimen-pre-corona}

\textbf{Autor:} T. Dufour\\
\textbf{Afiliación:} LPP, Sorbonne Université, CNRS, École
Polytechnique, Université Paris-Sud, Observatoire de Paris, Université
Paris-Saclay, PSL Research University, 4 Place Jussieu, 75252 Paris,
France\\
\textbf{Revista:} Compte-Rendus Mécanique de l'Académie des Sciences,
Volumen 354 (2026), pp.~117-139\\
\textbf{DOI:} \url{https://doi.org/10.5802/crmeca.346}

\begin{center}\rule{0.5\linewidth}{0.5pt}\end{center}

\subsection{Resumen}\label{resumen}

El electrovegetómetro de Pierre-Nicolas Bertholon del siglo XVIII fue
concebido para aprovechar la ``electricidad atmosférica'' para el
crecimiento de las plantas, pero sus capacidades físicas nunca han sido
cuantificadas dentro del contexto de la comprensión actual del sistema
eléctrico atmosférico terrestre. Este estudio aborda la falta de
evaluación cuantitativa de dicho dispositivo histórico de
``electrocultura'' y su plausible influencia en el entorno eléctrico
cercano al dosel. Su objetivo es reinterpretar el aparato de Bertholon
utilizando la electrodinámica atmosférica contemporánea, preguntándose
bajo qué condiciones realistas de tiempo despejado y de tormenta un
colector-distribuidor puramente pasivo podría generar campos y flujos
iónicos de significación física. Se ha desarrollado un modelo óhmico
cuasi-estacionario bidimensional en el que la atmósfera es una columna
resistente que transporta la corriente de conducción global, la
estructura metálica es un conductor flotante soportado por aislantes de
madera con fugas y se excluyen la carga espacial y la corona, de modo
que todos los resultados describen límites superiores pre-arranque. Las
simulaciones muestran que en tiempo despejado el punto superior simple y
la corona inferior de múltiples puntos del electrovegetómetro amplifican
el campo de fondo entre dos y tres órdenes de magnitud, pero solo dentro
de regiones milimétrico-centimétricas alrededor de las puntas y con
corrientes totales limitadas al rango de pA-nA·m⁻². Bajo forzamiento de
tormenta, los campos pico en la corona alcanzan 10⁵-10⁶ V·m⁻¹,
acercándose o excediendo los umbrales empíricos de arranque de corona,
mientras permanecen en gran medida insensibles a incertidumbres en el
ángulo de vértice o la geometría del colector siempre que esté presente
un mástil elevado. Estos resultados hacen que los informes de Bertholon
sobre ``aigrettes'' luminosas sean físicamente plausibles, pero sugieren
que cualquier impacto agronómico en tiempo despejado era sutil y
altamente localizado y que las reclamaciones modernas de
``electrocultura'' requieren estudios cuidadosos y acoplados de
electrostática y biología más allá del régimen pre-corona.

\textbf{Palabras clave:} Electrocultura, electrovegetómetro, agricultura
plasma, Bertholon

\begin{center}\rule{0.5\linewidth}{0.5pt}\end{center}

\section{Introducción}\label{introducciuxf3n}

\subsection{Un impulso interdisciplinario: desde el coloquio de
Montpellier hasta el presente
estudio}\label{un-impulso-interdisciplinario-desde-el-coloquio-de-montpellier-hasta-el-presente-estudio}

El presente estudio se inició como resultado del coloquio francés
\emph{«Homenaje a Pierre Bertholon, sabio y electricista de las Luces»},
celebrado en Montpellier en octubre de 2023 bajo el patrocinio de la
UNESCO. Este evento reunió a historiadores de la ciencia, físicos,
especialistas en electricidad atmosférica y académicos de la filosofía
natural del siglo XVIII, comunidades que rara vez interactúan aunque su
combinación posee un potencial interdisciplinario insospechado y
significativo. Las discusiones con los participantes durante y después
de la reunión revelaron un interés compartido en reevaluar el
electrovegetómetro de Bertholon utilizando las herramientas de la
electrodinámica atmosférica moderna, manteniendo fidelidad a su contexto
histórico y restricciones materiales. El presente artículo es, por lo
tanto, en parte producto de este encuentro interdisciplinario,
ilustrando el valor científico de tales coloquios para tender puentes
entre culturas de investigación distintas y estimular nuevas líneas de
investigación a través de fronteras disciplinarias.

\subsection{Encuadre histórico}\label{encuadre-histuxf3rico}

Para situar este renovado interés dentro de su contexto científico
original, es necesario revisitar los fundamentos históricos del trabajo
de Bertholon. En 1783, Pierre-Nicolas Bertholon de Saint-Lazare,
comúnmente conocido como Abbé Bertholon, publicó \emph{De l'électricité
des végétaux}, un tratado en el que afirmaba que la ``electricidad
atmosférica'' actúa continuamente sobre las plantas y podría
aprovecharse para mejorar la germinación, el crecimiento, la floración y
la fructificación {[}Bertholon, 1783{]}. Su propuesta práctica clave, el
electrovegetómetro (\emph{électro-végétomètre}), era un aparato pasivo
destinado a recolectar carga de la atmósfera y redistribuirla suavemente
sobre los cultivos a través de ``aigrettes lumineuses'': una expresión
francesa histórica que puede entenderse hoy como descargas de corona
luminosas tenues.

Aunque Bertholon interpretó el funcionamiento de su aparato a través de
la noción del siglo XVIII de ``fluido eléctrico'', el principio físico
al que apelaba puede expresarse claramente en términos electrodinámicos
modernos: conductores afilados sumergidos en un campo eléctrico
ambiental pueden actuar como fuentes (o sumideros) de carga espacial,
modificando las densidades iónicas locales y los flujos de aire débiles
y alterando así el entorno microfísico cerca de las hojas y el suelo. En
las descripciones de Bertholon, el electrovegetómetro consistía
esencialmente en un mástil vertical, aislado, que elevaba una estructura
de recolección hacia el campo eléctrico atmosférico. Una ruta conductora
llevaba la carga interceptada hacia abajo hasta un brazo distribuidor,
que terminaba en una corona de puntos finos suspendida sobre las hileras
de cultivos. Los puntos estaban destinados a liberar carga a través de
descargas no disruptivas, modificando sutilmente el gradiente de
potencial y el entorno de carga alrededor del dosel.

Bertholon describió dos configuraciones estrechamente relacionadas de
este electrovegetómetro:

\begin{itemize}
\tightlist
\item
  La Configuración I ``varilla de punto único, brazo articulado''
  (\textbf{Figura 1a}), combinaba un colector relativamente simple con
  un difusor mecánicamente
\end{itemize}

\begin{quote}
versátil. Un alto mástil de madera, cuidadosamente aislado del suelo por
alquitrán, resinas y elementos de vidrio, llevaba un punto de hierro
único en su cabeza, protegido por un embudo de hojalata para favorecer
la recolección ascendente. Una cadena suspendida conducía el ``fluido
eléctrico'' interceptado hasta un pequeño disco de hierro donde se
esperaba que la carga recolectada se acumulara antes de ser
redistribuida. Desde allí, un brazo metálico articulado, soportado
horizontalmente sobre caballetes aislantes, podía balancearse sobre
hileras seleccionadas de cultivos. En su extremo final, una corona de
puntos finos podía liberar carga como una corona difusa justo encima del
dosel. Una cadena de puesta a tierra removible permitía que todo el
aparato se neutralizara cuando el aire era juzgado ``sobrecargado''.
\end{quote}

\begin{itemize}
\tightlist
\item
  La Configuración II ``corona de múltiples puntos, brazo giratorio y
  extensible'' (\textbf{Figura 1b}) se caracteriza por una corona de
  puntos afilados montada en un cubo tratado con resina, aislado con
  vidrio,
\end{itemize}

\begin{quote}
para aumentar el área efectiva de recolección de \emph{fluido eléctrico}
manteniendo un alto aislamiento. Una palanca curvada y una cadena
conectaban este cubo a un brazo giratorio telescópico que podía barrer
un círculo completo y extenderse o retraerse para ajustar la huella
tratada. Como en la Configuración I, una corona terminal de puntos
estaba destinada a difundir carga suavemente sobre la vegetación.
\end{quote}

\begin{quote}
En ambas configuraciones, el aparato era enteramente pasivo: ni medía
voltaje ni corriente e inyectaba energía externa. Dado que conceptos
modernos como el potencial eléctrico, la corriente y las unidades
estandarizadas aún no estaban disponibles, Bertholon razonaba
cualitativamente con nociones de ``más'' o ``menos'' electricidad. Los
componentes individuales de cada configuración se detallan en la Tabla
1.
\end{quote}

\begin{quote}
\includegraphics{./images/image1.jpeg}
\end{quote}

\begin{quote}
\textbf{\emph{Figura 1. Electrovegetómetro de Bertholon, un aparato
pasivo destinado a canalizar la ``electricidad atmosférica'' hacia los
cultivos. (a) Configuración I, con un colector de varilla de punto único
(E) en la cabeza del mástil y un brazo articulado que distribuía la
carga a través de una corona de múltiples puntos (N). (b) Configuración
II, equipada con un colector de corona de múltiples puntos (E), un brazo
giratorio y extensible y una huella de tratamiento más grande que
termina en un distribuidor de corona de múltiples puntos (N). Un
dispositivo de neutralización (Z), mostrado en el recuadro de (b), podía
proporcionar seguridad y control de encendido/apagado. Adaptado de
{[}Bertholon, 1783{]}.}}
\end{quote}

\textbf{Tabla 1. Comparación entre el diseño y funcionamiento del
electrovegetómetro de Bertholon (Figura 1): Configuración I y
Configuración II.}

\begin{longtable}[]{@{}
  >{\raggedright\arraybackslash}p{(\columnwidth - 4\tabcolsep) * \real{0.1039}}
  >{\raggedright\arraybackslash}p{(\columnwidth - 4\tabcolsep) * \real{0.3961}}
  >{\raggedright\arraybackslash}p{(\columnwidth - 4\tabcolsep) * \real{0.5000}}@{}}
\toprule\noalign{}
\begin{minipage}[b]{\linewidth}\raggedright
Característica
\end{minipage} & \begin{minipage}[b]{\linewidth}\raggedright
Configuración I (varilla de punto único, brazo articulado)
\end{minipage} & \begin{minipage}[b]{\linewidth}\raggedright
Configuración II (corona de múltiples puntos, brazo giratorio y
extensible)
\end{minipage} \\
\midrule\noalign{}
\endhead
\bottomrule\noalign{}
\endlastfoot
\textbf{Propósito} & Captura pasiva → Conducción aislada → Difusión
suave sobre cultivos (sin medición) & Igual \\
\textbf{Mástil / Base} & Mástil de madera (A-B): sección enterrada
secada al fuego, alquitranada, envuelta en polvo de carbón/cemento,
luego base de albañilería; parte sobre el suelo pintada/con betún &
Mástil (A-B) con cilindro de cabeza saturado en resina (C) tratado con
alquitrán/brea/turpentina \\
\textbf{Aislamiento superior} & Tubo de vidrio grueso relleno de
mastique bituminoso (D); segundo aislante de vidrio aguas abajo (H) &
Interfaz rotativa aislada mediante inserción de resina (D) y manga de
vidrio/mastique (H) \\
\textbf{Colector (en la cabeza del mástil)} & Varilla de punto único (E)
montada en D; protegida por un embudo de hojalata (F) & Corona de
múltiples puntos (E) para maximizar la captura de carga \\
\textbf{Ruta de conducción (vertical)} & Cadena conductora (G)
suspendida de E, guiada a través de H para prevenir fugas & Palanca de
hierro curvada (C) que lleva la cadena (G) a la interfaz del brazo \\
\textbf{Regulación / buffer} & La cadena termina en un disco de hierro
(K) que actúa como pequeño ``condensador/regulador'' & Sin disco
separado; regulación implícita (enfoque en movilidad y alcance) \\
\textbf{Brazo horizontal} & Brazo conductor articulado (K-L-M-N) con
bisagras en L y Q, soportado sobre caballetes aislantes (O, P) con
cuerdas de seda tensadas & Brazo conductor telescópico (L-M) con
conductor telescópico (Q), soportado sobre caballete aislante (P) \\
\textbf{Difusor terminal} & Corona de puntos metálicos afilados (N, R),
que liberan corona no disruptiva & Corona de puntos metálicos
afilados \\
\textbf{Movilidad / cobertura} & El brazo se balancea en bisagras (L, Q)
para cubrir un sector (hilera/franja) & Rotación completa de 360°;
extender/retraer para variar el alcance \\
\textbf{Neutralización} & Cadena entre K y tierra para drenar carga &
Igual \\
\textbf{Seguridad} & ``Excitador en forma de C'' (marco de cobre/hierro
con mango de vidrio + cadena de tierra arrastrándose) & Igual \\
\end{longtable}

\subsection{De las ideas de la Ilustración a la electrodinámica
atmosférica
moderna}\label{de-las-ideas-de-la-ilustraciuxf3n-a-la-electrodinuxe1mica-atmosfuxe9rica-moderna}

A finales del siglo XVIII, varios aspectos esenciales de la electricidad
atmosférica ya eran reconocidos. Los experimentos de Stephen Gray y
Charles François de Cisternay Du Fay ya habían distinguido conductores
de aislantes e introducido la idea de dos estados eléctricos opuestos,
más tarde renombrados ``positivo'' y ``negativo'' {[}Bailey, 2001{]}.
Benjamin Franklin sintetizó estas ideas, definió la convención de
polaridad y argumentó, sobre la base de similitudes de laboratorio y
experimentos de campo, que el rayo era una descarga eléctrica de la
misma naturaleza que las chispas de las máquinas {[}Williams, 2009{]}.
Sus propuestas fueron confirmadas en Francia por Dalibard y Le Monnier
quienes utilizaron conductores de hierro elevados y bien aislados para
extraer chispas durante tormentas, confirmando así que las nubes
tormentosas portan carga eléctrica {[}Winkler et al., 1898{]}. Poco
después, las mediciones sistemáticas de electroscopio de John Canton en
Londres y Giambattista Beccaria en Turín mostraron que incluso el aire
sin nubes rara vez era neutro: descargas lentas ``silenciosas'',
variaciones diurnas del campo y brillos luminosos como el fuego de San
Elmo de los marineros revelaban que la atmósfera inferior estaba
persistentemente electrificada {[}Herbert, 2012{]}. Desde un punto de
vista moderno, el fuego de San Elmo se entendía como una descarga de
corona de baja corriente de conductores afilados en un campo ambiental
fuerte {[}McGinness, 2025{]}. Sin embargo, estos conocimientos del siglo
XVIII eran locales y cualitativos. Un marco cuantitativo y global, en el
que las tormentas, las regiones de tiempo despejado y la ionosfera están
acopladas a través de un solo circuito eléctrico atmosférico, surgió
solo a principios del siglo XX, particularmente

a través del trabajo de C. T. R. Wilson y las síntesis posteriores del
``circuito eléctrico atmosférico global'' {[}Aplin, 2008{]}.

La electrodinámica atmosférica moderna enmarca estos fenómenos en
términos del circuito eléctrico global (GEC). El sistema
Tierra-ionosfera se comporta como un capacitor esférico con fugas: las
tormentas y nubes electrificadas mantienen la ionosfera a
aproximadamente +250/+300 kV en relación con la superficie, mientras que
las regiones de ``tiempo despejado'' cierran el circuito a través de la
débil conductividad eléctrica del aire {[}Tinsley, 2007{]}, {[}Mareev,
2014{]}. En condiciones claras sin perturbaciones, la columna
atmosférica transporta una pequeña corriente de conducción descendente
casi estable, descrita localmente por la ley de Ohm:

\begin{quote}
𝐽\textsubscript{𝑧}(𝑧) = 𝜎(𝑧) 𝐸\textsubscript{𝑧}(𝑧) \{1\}
\end{quote}

donde σ(z) es la conductividad eléctrica y Ez(z) el campo eléctrico
vertical (tomado positivo hacia arriba). Cerca de la superficie, los
valores típicos son E0,atm ≈ −100 a −150 V.m⁻¹ y J0,atm ≈ −1 a −3 pA.m⁻²
{[}Guha, 2010{]}. La conductividad aumenta muchos órdenes de magnitud
con la altitud (desde aproximadamente 10⁻¹⁴ S.m⁻¹ cerca del suelo hasta
10⁻⁷ S.m⁻¹ en la termosfera inferior), de modo que
\textbar 𝐸\textsubscript{𝑧}\textbar{} disminuye correspondientemente y
el potencial V(z) aumenta desde 0 V en el suelo hacia el nivel
ionosférico de unos pocos ×10⁵ V {[}Kudintseva, 2016{]}.

Bajo condiciones perturbadas o tormentosas, esta estructura vertical
simple se modifica fuertemente. La separación de carga en las nubes
tormentosas produce campos cuasi-estáticos grandes en la superficie
(hasta varios kV.m⁻¹), aumenta la producción de iones y altera el perfil
de conductividad {[}Kundt, 1999{]}. La densidad de corriente vertical
puede aumentar al rango de nA.m⁻² y se desarrollan campos fuertes y
localmente variables cerca de objetos afilados, creando condiciones
favorables para descargas de corona y leaders {[}Standler, 1979{]}.
Estos estados perturbados forman la rama de tormenta del circuito global
y son cruciales para interpretar cómo un conductor que sobresale, como
el electrovegetómetro de Bertholon, podría experimentar tanto un
forzamiento modesto de tiempo despejado como un forzamiento episódico y
mucho más fuerte de tiempo de tormenta.

\textbf{Tabla 2. Regímenes superficiales idealizados utilizados como
condiciones de contorno atmosféricas.}

\begin{longtable}[]{@{}
  >{\raggedright\arraybackslash}p{(\columnwidth - 4\tabcolsep) * \real{0.1642}}
  >{\raggedright\arraybackslash}p{(\columnwidth - 4\tabcolsep) * \real{0.4776}}
  >{\raggedright\arraybackslash}p{(\columnwidth - 4\tabcolsep) * \real{0.3582}}@{}}
\toprule\noalign{}
\begin{minipage}[b]{\linewidth}\raggedright
Parámetro
\end{minipage} & \begin{minipage}[b]{\linewidth}\raggedright
Condiciones de tiempo despejado
\end{minipage} & \begin{minipage}[b]{\linewidth}\raggedright
Condiciones de tormenta
\end{minipage} \\
\midrule\noalign{}
\endhead
\bottomrule\noalign{}
\endlastfoot
\textbf{Dirección del campo eléctrico} & Descendente (negativo) &
Ascendente (positivo) \\
\textbf{Campo eléctrico vertical E₀,atm} & −100 → −150 V·m⁻¹ & +1 → +10
kV·m⁻¹ \\
\textbf{Potencial eléctrico V(z)} & Aumenta con la altitud (dV/dz =
−E\_z \textgreater{} 0) & Disminuye con la altitud (dV/dz = −E\_z
\textless{} 0) \\
\textbf{V(1m)} & +100 V → +150 V & −1 kV → −10 kV \\
\textbf{V(10m)} & +1000 V → +1500 V & −10 kV → −100 kV \\
\textbf{Conductividad en superficie σ(0)} & (1-5) × 10⁻¹⁴ S·m⁻¹ & (1-10)
× 10⁻¹² S·m⁻¹ \\
\textbf{Variación de σ con la altura (0-20 m)} & Aumenta gradualmente
(≈10×) debido a la ionización por rayos cósmicos y radiactividad &
Aumenta bruscamente (≈100×) cerca del suelo debido a la ionización
inducida por campo y carga espacial \\
\textbf{Densidad de corriente de conducción J} & ≈ −2 pA·m⁻²
(descendente, negativo) & +0,1 → +10 nA·m⁻² (ascendente, positivo) \\
\textbf{Referencia de potencial superficial V(0)} & 0 V (por convención)
& 0 V (por convención) \\
\textbf{Campo eléctrico de ruptura (Ebd)} & ≈ 3,0 MV·m⁻¹ & 2,5-2,8
MV·m⁻¹ \\
\end{longtable}

Debido a su altura limitada (\textbf{Figura 1}), el aparato de Bertholon
solo puede interactuar con los primeros metros de la atmósfera, donde la
imagen idealizada unidimensional del GEC se rompe y los procesos locales
dominan. La topografía, la carga de aerosoles, la vegetación, la
humedad, la conductividad superficial y la proximidad a la carga de
tormenta influyen todos en el campo eléctrico local. Para un
colector-distribuidor pasivo como el electrovegetómetro, esta capa
cercana a la superficie controla tanto la magnitud de la corriente
interceptada como la probabilidad de iniciar descargas en las puntas.
Para mayor claridad, se pueden distinguir dos regímenes idealizados,
como se resume en la Tabla 2:

\begin{itemize}
\item
  Condiciones de tiempo despejado (``soleado''), que representan estados
  atmosféricos sin perturbaciones con campos débiles y corrientes de
  conducción muy pequeñas;
\item
  Condiciones de tiempo tormentoso, que representan períodos bajo nubes
  electrificadas o cerca de celdas tormentosas activas, donde los campos
  y las corrientes pueden aumentar por órdenes de magnitud.
\end{itemize}

\subsection{Brecha de conocimiento y preguntas de
investigación}\label{brecha-de-conocimiento-y-preguntas-de-investigaciuxf3n}

Los relatos históricos describen el electrovegetómetro de Bertholon como
un aparato puramente pasivo: la corriente de conducción atmosférica
ambiental

era recolectada en una punta superior afilada y redistribuida a través
de un conductor por una corona inferior de puntos, sin ninguna fuente de
energía externa o medio para regular o medir el estado eléctrico. Lo que
sigue siendo incierto es si una configuración pasiva de este tipo
podría, incluso en principio, generar flujos iónicos o de carga lo
suficientemente grandes como para tener significación física a la altura
del dosel. Ningún estudio anterior ha incrustado explícitamente el
electrovegetómetro en el marco del circuito eléctrico atmosférico
global, con campos, corrientes y perfiles de conductividad realistas de
tiempo despejado y de tormenta, ni ha cuantificado hasta dónde desde las
puntas su influencia en el microentorno eléctrico de los cultivos puede
extenderse.

Una segunda brecha concierne la ausencia de una distinción sistemática
entre los regímenes de tiempo despejado y de tormenta. La electricidad
atmosférica moderna reconoce que las condiciones típicas de tiempo
despejado (por ejemplo, 𝐸\textsubscript{0,𝑎𝑡𝑚} ≈ −120 𝑘𝑉.
𝑚\textsuperscript{−1}, 𝐽\textsubscript{0,𝑎𝑡𝑚} ≈ −2 𝑝𝐴.
𝑚\textsuperscript{−2}) se asocian con corrientes de conducción débiles y
cuasi-estacionarias, mientras que las condiciones de tiempo de tormenta
pueden soportar campos eléctricos mucho más fuertes y estructurados. Sin
embargo, las discusiones historiográficas y experimentales existentes
del electrovegetómetro no cuantifican cómo estos regímenes impulsarían
diferencialmente una estructura metálica flotante, ni estiman los
límites superiores pre-corona de los campos y corrientes alcanzables en
el colector superior y la corona inferior. Esto dificulta juzgar si los
informes históricos de ``aigrettes'' luminosas corresponden a campos que
plausiblemente exceden los umbrales de arranque de corona difusa.

Finalmente, las consecuencias electrodinámicas de la incertidumbre
geométrica siguen sin explorarse. Los grabados históricos restringen la
arquitectura general del electrovegetómetro pero no la altura precisa
del mástil, la longitud del brazo, la separación de las puntas ni los
ángulos de vértice. Por lo tanto, se desconoce cuán fuertemente la
corriente de columna de fondo puede concentrarse en la capa del dosel
por una corona realista, sobre qué distancia sobre el suelo persisten el
campo eléctrico y la densidad de corriente mejorados (por ejemplo,
dentro de los primeros 0,2-0,5 m) y cuán sensibles son estas cantidades
a la geometría del colector y del mástil.

El presente trabajo aborda estas brechas desarrollando un modelo
simplificado de conducción bidimensional del sistema aire-dispositivo,
en el que el electrovegetómetro se trata como una estructura metálica
flotante polarizada por el circuito eléctrico global. Dentro de este
marco óhmico pre-corona, calculamos el potencial, el campo eléctrico y
la densidad de corriente alrededor de reconstrucciones idealizadas del
dispositivo de Bertholon bajo forzamiento de tiempo despejado y de
tormenta, para varias combinaciones plausibles de altura del mástil,
longitud del brazo y separación de puntas. Las simulaciones cuantifican
(i) los campos locales máximos y los factores de amplificación en
colectores de punto único y de múltiples puntos, (ii) la concentración
de la corriente de columna de fondo en la capa del dosel y los
gradientes verticales y horizontales resultantes de \textbar 𝐸\textbar{}
y \textbar 𝐽\textbar{} y (iii) las condiciones atmosféricas y
geométricas bajo las cuales los campos calculados se acercan o exceden
los criterios empíricos de arranque de corona difusa. Dado que no se
modela la acumulación de carga espacial ni la retroalimentación de
ionización, estos resultados deben interpretarse como límites superiores
controlados pre-arranque en lugar de predicciones del comportamiento de
descarga.

Estas consideraciones llevan a las siguientes preguntas de
investigación:

\begin{quote}
\begin{itemize}
\tightlist
\item
  Bajo forzamiento realista de tiempo despejado, qué potenciales,
  magnitudes de campo eléctrico y densidades de corriente se establecen
  alrededor del colector superior y la corona inferior y
\end{itemize}
\end{quote}

\begin{quote}
sobre qué escalas de longitud decaen estas perturbaciones sobre el
dosel?
\end{quote}

\begin{quote}
ecuación \{5\}, donde 𝜎\textsubscript{0}
\end{quote}

\begin{quote}
= {[}𝐽0,atm{]} está determinada por el 𝐸0,atm impuesto
\end{quote}

\begin{itemize}
\item
  Bajo campos de fondo mejorados que representan condiciones
  pre-tormenta o afectadas por tormenta, ¿puede un electrovegetómetro
  puramente pasivo alcanzar niveles de campo eléctrico en sus puntas que
  se acerquen a los umbrales de arranque de corona difusa, haciendo que
  las ``aigrettes'' luminosas históricas sean físicamente plausibles?
\item
  ¿Cómo dependen la amplificación del campo, la concentración de
  corriente en la capa del dosel y la proximidad a los umbrales de
  arranque de corona de la altura del mástil, la longitud del brazo, la
  separación de puntas y el diseño del colector dentro del rango
  consistente con las descripciones históricas?
\end{itemize}

\section{Modelo}\label{modelo}

Todos los cálculos numéricos se realizaron utilizando GNU Octave
(versión 10.3.0), empleando scripts personalizados para la
discretización de diferencias finitas, el esquema de relajación y el
post-procesamiento.

\subsection{Supuestos físicos y
alcance}\label{supuestos-fuxedsicos-y-alcance}

El entorno electrostático alrededor del electrovegetómetro de Bertholon
se modela como una atmósfera cuasi-estacionaria y óhmica que transporta
la corriente de conducción global del clima. Los supuestos clave son los
siguientes:

campo eléctrico atmosférico superficial 𝐸\textsubscript{0,atm} y la

densidad de corriente de columna correspondiente 𝐽\textsubscript{0,atm}.
\textbf{La Figura 2a} muestra la estructura de isopotenciales resultante
para valores de 𝐸\textsubscript{0,atm} que van desde −0,4 𝑘𝑉.
𝑚\textsuperscript{−1} hasta 10 𝑘𝑉. 𝑚\textsuperscript{−1}. Aparecen dos
regímenes atmosféricos físicamente distintos en este rango y, por lo
tanto, se consideran en las simulaciones:

\begin{itemize}
\tightlist
\item
  Condiciones de tiempo despejado, caracterizadas por
  𝐸\textsubscript{0,atm} ≈
\end{itemize}

\begin{quote}
−120𝑉. 𝑚\textsuperscript{−1} y 𝐽\textsubscript{0,atm} ≈ −2 𝑝𝐴.
𝑚\textsuperscript{−2}, produciendo gradientes de potencial débiles e
isopotenciales suavemente curvados (lado izquierdo de la \textbf{Figura
2a}).
\end{quote}

\begin{itemize}
\tightlist
\item
  Condiciones de tiempo tormentoso, asociadas con campos eléctricos
  positivos mucho más fuertes, aquí representados por
  𝐸\textsubscript{0,atm} =
\end{itemize}

\begin{quote}
+5𝑘𝑉. 𝑚\textsuperscript{−1} y 𝐽\textsubscript{0,atm} = 5𝑛𝐴.
𝑚\textsuperscript{−2}, que producen caídas de potencial pronunciadas con
la altitud e isopotenciales densamente empaquetados (lado derecho de la
\textbf{Figura 2a}).
\end{quote}

\begin{itemize}
\tightlist
\item
  Se adopta una altura de escala de conductividad de 𝐻\textsubscript{𝜎}
  = 6𝑘𝑚. En una columna atmosférica verticalmente uniforme, la
  conservación de corriente \{6\} conduce entonces al perfil de campo
  eléctrico de la ecuación \{7\}. A través de la ley de Gauss, la
  densidad de carga espacial de diagnóstico correspondiente viene dada
  por la ecuación \{8\}.
\end{itemize}

∇⃗→ ⋅ {[}𝜎(𝑥, 𝑧) ∇⃗→𝑉(𝑥, 𝑧){]} = 0 \{2\}

𝐸⃗→ = −∇⃗→𝑉 \{3\}

𝐽→ = 𝜎𝐸⃗→ \{4\}

\begin{itemize}
\tightlist
\item
  La atmósfera se comporta como un medio lineal y resistivo
  caracterizado por una conductividad prescrita σ(x, z), dominada
\end{itemize}

𝜎air

(𝑧) = 𝜎\textsubscript{0}

\begin{quote}
{[}𝑧{]}
\end{quote}

exp ( ) \{5\}

\begin{quote}
𝐻\textsubscript{𝜎}
\end{quote}

\begin{quote}
por su dependencia de altitud σ(z).
\end{quote}

\begin{itemize}
\tightlist
\item
  La corriente de conducción vertical se conserva aproximadamente
\end{itemize}

\begin{quote}
𝐽\textsubscript{𝑧}(𝑧) = 𝜎(𝑧). 𝐸\textsubscript{𝑧}(𝑧) =
𝐽\textsubscript{0,𝑎𝑡𝑚} \{6\}
\end{quote}

\begin{quote}
(∇⃗→. →J = 0) en ausencia de fuentes o sumideros. La carga espacial ρ y
la retroalimentación de ionización no se resuelven explícitamente; en
cambio, su efecto neto sobre la conductividad está codificado en σ(z).
\end{quote}

\begin{quote}
𝐸\textsubscript{𝑧}(𝑧) = −𝐸\textsubscript{0,𝑎𝑡𝑚}
\end{quote}

𝑑𝐸\textsubscript{𝑧}

\begin{quote}
𝑧
\end{quote}

exp (− ) \{7\}

𝐻

\begin{quote}
𝜎
\end{quote}

\begin{quote}
𝐸\textsubscript{𝑧}(𝑧)
\end{quote}

\begin{itemize}
\tightlist
\item
  Las partes metálicas del dispositivo se tratan como conductores
  perfectos, cada región metálica conectada siendo
\end{itemize}

\begin{quote}
𝜌(𝑧) = 𝜀\textsubscript{0} 𝑑𝑧 = −𝜀\textsubscript{0}
\end{quote}

\begin{quote}
𝐻\textsubscript{𝜎}
\end{quote}

\begin{quote}
\{8\}
\end{quote}

\begin{quote}
equipotencial. La estructura es eléctricamente flotante: no se impone
voltaje externo y su potencial se ajusta de manera autoconsistente bajo
el forzamiento del circuito global impuesto.
\end{quote}

\begin{itemize}
\tightlist
\item
  El suelo se representa como un límite conductor y equipotencial. El
  dosel de cultivos no se resuelve explícitamente (los campos se evalúan
  en el aire sobre un suelo plano), por lo que las distorsiones de campo
  a pequeña escala
\end{itemize}

\begin{quote}
La densidad de carga espacial 𝜌(𝑧) no se resuelve explícitamente. En
cambio, σ(z) se prescribe y el solucionador permanece puramente óhmico y
pre-corona. Para casos sustitutos de tipo tormentosa, se considera la
misma forma aunque las magnitudes de
\textbar 𝐸\textsubscript{0,𝑎𝑡𝑚}\textbar,
\textbar 𝐽\textsubscript{0,𝑎𝑡𝑚}\textbar{} y 𝜎\textsubscript{0} se
aumentan para igualar los valores idealizados en el régimen ``tipo
tormenta''. Las regiones
\end{quote}

\begin{quote}
sólidas son aire con 𝜎\textsubscript{air}(𝑧) dado por la ecuación \{5\},
soportes de madera
\end{quote}

\begin{quote}
por plantas individuales se descuidan.
\end{quote}

\begin{itemize}
\tightlist
\item
  El problema se resuelve en dos dimensiones {[}Cartesiano /
\end{itemize}

\begin{quote}
con 𝜎\textsubscript{wood}
\end{quote}

≈ 10\textsuperscript{−15} 𝑆. 𝑚\textsuperscript{−1} (fuertemente aislante
pero no perfectamente)

\begin{quote}
7 −1
\end{quote}

\begin{quote}
axisimétrico{]} geométrico, asumiendo invarianza en la dirección
perpendicular al plano.
\end{quote}

\begin{itemize}
\tightlist
\item
  Los soportes de madera se representan como conductores débiles (σwood
  muy baja) que pueden tener fugas de corriente pero son altamente
  resistivos comparados con el metal.
\end{itemize}

\subsection{Ecuaciones gobernantes y perfil de
conductividad}\label{ecuaciones-gobernantes-y-perfil-de-conductividad}

El potencial eléctrico 𝑉(𝑥, 𝑧) satisface la ecuación de conducción
estacionaria \{2\} en todos los materiales, mientras que el campo
eléctrico y la densidad de corriente se derivan de las ecuaciones \{3\}
y \{4\}. La atmósfera se describe utilizando el perfil de conductividad
dependiente de la altitud de

y el conductor metálico: 𝜎\textsubscript{metal} = 1 × 10 𝑆. 𝑚. Debido a
esta

fuerte discrepancia de conductividad, la estructura metálica se trata
como un solo conductor flotante con un potencial interno efectivamente
uniforme.

\subsection{Geometría y dominio
computacional}\label{geometruxeda-y-dominio-computacional}

El problema se resuelve en dos dimensiones sobre 𝑥 ∈ {[}−0,4, 2,0{]} 𝑚 y
𝑧 ∈ {[}0,0, 4,0{]} 𝑚 en una cuadrícula uniforme de 𝑁\textsubscript{𝑥} =
6000, 𝑁\textsubscript{𝑧} = 10 000 nodos (Δx = Δz = 400 µm). Esta
formulación 2D representa una estructura infinitamente larga en la
dirección perpendicular al plano; los resultados son, por lo tanto,
secciones transversales con profundidad unitaria.

El aparato tipo Bertholon se modela como se muestra en la \textbf{Figura
2b}. Un mástil vertical de madera, de 2,0 m de alto y 0,15 m de ancho,
forma el soporte central. Un elemento horizontal corto de madera se
coloca a

𝑧 = 1,5 m y dos pies de madera, cada uno de 0,5 m de alto y 1,5 cm de
grosor, proporcionan soporte a nivel del suelo debajo del conductor. El
conductor metálico se representa por una estructura en forma de L con
puntas triangulares. En la parte superior, una punta triangular de 0,30
m de alto y 0,02 m de ancho ubicada en 𝑧 = 2,2 m sirve como el colector
superior de punto único, con un ángulo de vértice de 4°. Desde la base
de esta punta, un segmento metálico vertical se extiende 1,70 m hacia
abajo con un grosor de

0,02 m. Este segmento se conecta a un brazo horizontal de 1,50 m de
largo

posicionado aproximadamente a 0,5 m sobre el suelo. Al final de este
brazo, una corona terminal compuesta por cinco puntas triangulares (cada
una de 0,30 m de alto y 0,02 m en la base) se dispone con orientaciones
angulares de −40°, −20°, 0°, +20° y +40° con respecto a la vertical.
Esta geometría estilizada corresponde a radios de vértice agudos del
orden de 1 mm.

\begin{quote}
\includegraphics{./images/image2.jpeg}
\end{quote}

\begin{quote}
\textbf{\emph{Figura 2. (a) Potencial eléctrico}} 𝑽(𝑬𝒛, 𝒛)
\textbf{\emph{en la atmósfera sobre el suelo, mostrado como contornos de
colores (en kV) y líneas de isopotencial negras. El eje horizontal
representa el campo eléctrico vertical impuesto}} 𝑬𝒛\textbf{\emph{, que
va desde valores de tiempo despejado (}}𝑬𝒛 \textless{} 𝟎\textbf{\emph{)
hasta valores de tiempo de tormenta (}}𝑬𝒛 \textgreater{}
𝟎\textbf{\emph{). Las dos líneas verticales blancas corresponden a las
condiciones atmosféricas de tiempo despejado y tormentoso de este
estudio. (b) Modelo geométrico del electrovegetómetro de Bertholon
implementado en la configuración I, que incluye el mástil aislante de
madera y los soportes (marrón), la estructura conductora en forma de L
(rojo) y la corona de múltiples puntos. El dominio computacional
representa una sección transversal vertical de 2 m × 4 m de la atmósfera
circundante.}}
\end{quote}

\subsection{Condiciones de contorno e implementación
numérica}\label{condiciones-de-contorno-e-implementaciuxf3n-numuxe9rica}

\begin{quote}
implementación
\end{quote}

Los límites superior e inferior aplican la densidad de corriente
vertical impuesta 𝐽\textsubscript{0,𝑎𝑡𝑚} siguiendo las ecuaciones \{9\}
y \{10\} mientras que los límites laterales son simétricos, como se
expresa en la ecuación \{11\}.

La ecuación gobernante \{2\} se discretiza en un esquema de

diferencias finitas de cinco puntos bidimensional. Para asegurar la
continuidad de la densidad de corriente a través de cambios abruptos de
conductividad, la conductividad en cada cara de celda se calcula con la
ecuación \{12\} que proporciona la media armónica entre celdas
adyacentes. Este esquema de promediado preserva la conservación de flujo
y la

𝜕𝑉

𝜕𝑧

𝑧=0

\begin{quote}
= − 𝐽0,𝑎𝑡𝑚 \{9\}
\end{quote}

\begin{quote}
𝜎(0)
\end{quote}

\begin{quote}
estabilidad en regiones donde materiales de conductividades muy
diferentes se encuentran (por ejemplo, interfaces aire/madera o
aire/metal). El
\end{quote}

\begin{quote}
sistema discreto resultante se resuelve por relajación de Jacobi, en el
que el
\end{quote}

{[}𝜕𝑉{]}

𝜕𝑧

𝑧=𝐻

\begin{quote}
{[}𝐽0,𝑎𝑡𝑚{]}
\end{quote}

\begin{quote}
= − 𝜎(𝐻)
\end{quote}

\{10\}

\begin{quote}
potencial en el nodo (𝑖, 𝑗) se actualiza iterativamente según la
ecuación
\end{quote}

\begin{quote}
\{13\} con actualizaciones de contorno y aplicación de calibre después
de cada
\end{quote}

𝜕𝑉

𝜕𝑥

izquierda/derecha

= 0 \{11\}

\begin{quote}
barrido.
\end{quote}

\begin{quote}
{[}2𝜎\textsubscript{1}𝜎\textsubscript{2}{]}
\end{quote}

\begin{quote}
𝜎 =
\end{quote}

\begin{quote}
\{12\}
\end{quote}

Un nodo de aire se fija en 𝑉 = 0 para establecer el potencial de
referencia. Esta elección aplica la corriente de columna pero no impone
una Tierra perfectamente equipotencial en z = 0; los potenciales
cercanos al suelo pueden

\begin{quote}
(𝑛+1)
\end{quote}

𝑉

\begin{quote}
𝑖,𝑗
\end{quote}

\begin{quote}
cara
\end{quote}

𝜎\textsubscript{1} + 𝜎\textsubscript{2} + 𝜀

por lo tanto variar horizontalmente, una aproximación aceptable aquí
pero no exacta.

\begin{quote}
𝜎𝑒𝑉𝑖,𝑗+1 + 𝜎𝑤𝑉𝑖,𝑗−1 + 𝜎𝑛𝑉𝑖+1,𝑗 + 𝜎𝑠𝑉𝑖−1,𝑗
\end{quote}

\begin{quote}
= 𝜎 + 𝜎 + 𝜎 + 𝜎
\end{quote}

𝑒 𝑤 𝑛 𝑠

\{13\}

La inicialización utiliza el perfil analítico de tiempo despejado en la
ecuación \{14\} y las regiones metálicas se siembran con el potencial
promedio de las celdas de aire vecinas. Las iteraciones se detienen una
vez que se satisface la ecuación \{15\} o cuando se alcanzan 2000
iteraciones. Este criterio produce una buena concordancia con la columna
analítica 1D y la conservación de corriente (el flujo superior vs
inferior difiere en \textless0,01\%), pero los campos pico siguen siendo
algo sensibles a la resolución de la cuadrícula y la geometría de las
puntas. Deben interpretarse como valores pre-corona aproximados en lugar
de cantidades completamente convergidas. Las iteraciones continúan hasta
que el cambio máximo en V entre dos pasos sucesivos cae por debajo de
una tolerancia prescrita, asegurando la convergencia del campo de
potencial electrostático.

\section{Resultados}\label{resultados}

\subsection{Condiciones de tiempo despejado
(``soleado'')}\label{condiciones-de-tiempo-despejado-soleado}

Aquí examinamos la Configuración I del electrovegetómetro
(\textbf{Figura 1a}), que consiste en una estructura conductora en forma
de L que combina un colector de punto único en la extremidad superior
con una corona de cinco puntos cerca del nivel del suelo. La respuesta
de esta geometría se evalúa bajo forzamiento de tiempo despejado,
caracterizado por una corriente de conducción descendente impuesta
𝐽\textsubscript{0,𝑎𝑡𝑚} ≈ −2 𝑝𝐴. 𝑚\textsuperscript{−2} y un campo
eléctrico de fondo 𝐸\textsubscript{0,𝑎𝑡𝑚} ≈ −120𝑉. 𝑚\textsuperscript{−1}
aplicado en el límite inferior.

𝑉(0)

(𝑧) = 𝐸\textsubscript{0,atm}𝐻\textsubscript{𝜎} (1 − 𝑒

\begin{quote}
{[}𝑧{]}
\end{quote}

\begin{quote}
𝐻𝜎) \{14\}
\end{quote}

−

\begin{quote}
El potencial en estado estacionario, la magnitud del campo eléctrico y
la magnitud de la densidad de corriente en las cercanías del colector
superior de punto único
\end{quote}

\begin{quote}
Max\textbar 𝑉\textsuperscript{(𝑛+1)} − 𝑉\textsuperscript{(𝑛)}\textbar{}
\textless{} 10\textsuperscript{−3} \{15\}
\end{quote}

El campo eléctrico y la densidad de corriente se calculan después de la
convergencia utilizando diferencias finitas centradas. Los componentes
horizontal y vertical del campo eléctrico se derivan directamente de los
gradientes de potencial (Ecuaciones \{16\}-\{17\}) y la magnitud del
campo se obtiene de la ecuación \{18\}. Los componentes de la densidad
de corriente se derivan entonces de la ley de Ohm (Ecuaciones
\{19\}-\{20\}) usando la conductividad local y la magnitud total de la
densidad de corriente se calcula a partir de la Ecuación \{21\}. Estas
cantidades se visualizan en unidades de kV.m⁻¹ para el campo eléctrico y
pA.m⁻² o nA.m⁻² para la densidad de corriente, con regiones sólidas
enmascaradas.

\begin{quote}
𝐸 = − 𝜕𝑉 \{16\}
\end{quote}

\begin{quote}
\textsuperscript{𝑥} 𝜕𝑥
\end{quote}

se muestran en la \textbf{Figura 3}. La \textbf{Figura 3a} ilustra que
el potencial atmosférico varía desde aproximadamente 260 hasta 305 V
sobre la extensión vertical del dominio. Cerca del punto metálico, el
potencial cae bruscamente a ≈260-265 V, mientras que el aire circundante
a la misma altitud se encuentra alrededor de 290-300 V. Esta depresión
lateral de 30-40 V está firmemente confinada a una región de solo unos
pocos mm de ancho, formando un embudo de potencial estrecho alineado con
la geometría cónica de la varilla. Dado que el conductor es
eléctricamente flotante, su potencial (≈240 V en la base del punto)
surge de manera autoconsistente como el valor requerido para satisfacer
la corriente de columna impuesta. Fuera de la región inmediata de la
punta, los contornos de isopotencial relajan de nuevo hacia su espaciado
casi uniforme de tiempo despejado.

La \textbf{Figura 3b} revela la amplificación resultante del campo
eléctrico. El gradiente asociado con la caída de 30-40 V sobre
distancias sub-centimétricas produce campos locales que exceden 20 𝑘𝑉.
𝑚\textsuperscript{−1} con la simulación

𝐸\textsubscript{𝑧}

\begin{quote}
𝜕𝑉
\end{quote}

\begin{quote}
= − 𝜕𝑧
\end{quote}

\{17\}

\begin{quote}
dando un valor pico de 𝐸\textsubscript{max} = 20,9 𝑘𝑉.
𝑚\textsuperscript{−1}. Esto corresponde a un factor de amplificación de
aproximadamente:
\end{quote}

\begin{quote}
\textbar 𝐸\textbar{} = √𝐸\textsuperscript{2} + 𝐸\textsuperscript{2}
\{18\}
\end{quote}

𝐸max

2,1 × 10\textsuperscript{4}

\begin{quote}
𝑥 𝑧
\end{quote}

\begin{quote}
𝜕𝑉
\end{quote}

𝐽 = −𝜎.

\{19\}

\textbar 𝐸0,𝑎𝑡𝑚\textbar{}

≈ 1,2 × 10\textsuperscript{2} ∼ 175

\begin{quote}
\textsuperscript{𝑥} 𝜕𝑥
\end{quote}

\begin{quote}
𝜕𝑉
\end{quote}

\begin{quote}
𝐽\textsubscript{𝑧} = −𝜎. 𝜕𝑧 \{20\}
\end{quote}

𝑥 𝑧

\begin{quote}
\textbar 𝐽\textbar{} = √𝐽\textsuperscript{2} + 𝐽\textsuperscript{2}
\{21\}
\end{quote}

\begin{quote}
Este valor se debe puramente a la concentración geométrica en el
vértice.
\end{quote}

\begin{quote}
El campo elevado decae rápidamente: dentro de 1 cm de la punta cae al
rango de kV/m y dentro de unos pocos cm regresa a los valores
ambientales de tiempo despejado.
\end{quote}

\begin{quote}
La \textbf{Figura 3c} muestra la magnitud de la densidad de corriente
asociada. Las corrientes más fuertes están confinadas en un estrecho
chorro axial que emerge de la varilla de punto único, con un pico de
𝐽\textsubscript{max} = 1,36 𝑛𝐴. 𝑚\textsuperscript{−2}. Este valor es
consistente con la ley de Ohm 𝐽 = 𝜎𝐸, usando la conductividad
superficial de tiempo despejado local 𝜎\textsubscript{0} ≈ 1,7 ×
10\textsuperscript{−14} 𝑆. 𝑚\textsuperscript{−1}. Aunque pequeña en
términos absolutos, una corriente de este tipo corresponde a un flujo
microscópico de:
\end{quote}

\begin{quote}
𝑁̇ ≈ {[}𝐽max{]} ≈ 8 × 10\textsuperscript{9}
cargas.s\textsuperscript{−1}.pm\textsuperscript{−2},
\end{quote}

𝑒

\begin{quote}
mostrando que incluso las corrientes de nanoamperio involucran miles de
millones de portadores de carga por segundo. Como se espera en el
régimen óhmico pre-corona, el patrón espacial de \textbar 𝐽\textbar{}
refleja el de \textbar 𝐸\textbar{} y se fusiona suavemente con la
corriente atmosférica de fondo 𝐽\textsubscript{0,𝑎𝑡𝑚} ≈ −2 𝑝𝐴.
𝑚\textsuperscript{−2} sobre escalas de cm.
\end{quote}

\begin{quote}

\includegraphics{./images/image3.jpeg}
\end{quote}

\begin{quote}
\textbf{\emph{Figura 3. Respuesta pre-corona del colector superior de
varilla de punto único (ángulo de vértice}} θ \textbf{\emph{= 4°) bajo
forzamiento de tiempo despejado (}}𝑬𝟎,𝒂𝒕𝒎 ≈ −𝟏𝟐𝟎𝑽.
𝒎\textsuperscript{−𝟏}\textbf{\emph{,}} 𝑱𝟎,𝒂𝒕𝒎 ≈
\end{quote}

\begin{quote}
−𝟐𝒑𝑨. 𝒎\textsuperscript{−𝟐}\textbf{\emph{). (a) Potencial V}}
𝒔\textbf{\emph{mostrando fuerte compresión de isopotenciales cerca del
vértice metálico flotante. (b) Magnitud del campo eléctrico}}
\textbar 𝑬\textbar{}\textbf{\emph{con amplificación geométrica hasta}}
𝑬𝒎𝒂𝒙 = 𝟐𝟎. 𝟗 𝒌𝑽. 𝒎\textsuperscript{−𝟏}\textbf{\emph{. (c) Magnitud de la
densidad de corriente}} \textbar 𝑱\textbar{} \textbf{\emph{con pico en}}
𝑱𝒎𝒂𝒙 = 𝟏. 𝟑𝟔 𝒏𝑨. 𝒎\textsuperscript{−𝟐} \textbf{\emph{y relajándose al
fondo dentro de centímetros.}}
\end{quote}

\includegraphics{./images/image4.jpeg}

\begin{quote}
\textbf{\emph{Figura 4. Respuesta pre-corona de la corona inferior de
cinco puntos bajo forzamiento de tiempo despejado (}}𝑬𝟎,𝒂𝒕𝒎 ≈ −𝟏𝟐𝟎𝑽.
𝒎\textsuperscript{−𝟏}\textbf{\emph{,}} 𝑱𝟎,𝒂𝒕𝒎 ≈ −𝟐𝒑𝑨.
𝒎\textsuperscript{−𝟐}\textbf{\emph{) con ángulo de vértice}} θ
\end{quote}

\begin{quote}
\textbf{\emph{= 4°. (a,b) Potencial V mostrando depresiones locales
alrededor de cada vértice y un embudo de potencial estrecho y profundo
en la punta central. (c,d) Magnitud del campo eléctrico}}
\end{quote}

\begin{quote}
\textbar 𝑬\textbar{} \textbf{\emph{formando cinco lóbulos de alto campo
interactuantes, con un máximo localizado}} 𝑬𝒎𝒂𝒙 = 𝟖𝟎. 𝟗 𝒌𝑽.
𝒎\textsuperscript{−𝟏}\textbf{\emph{. (e,f) Magnitud de la densidad de
corriente}} \textbar 𝑱\textbar{} \textbf{\emph{revelando cinco canales
de conducción enfocados y un pico}} 𝑱𝒎𝒂𝒙 = 𝟏. 𝟑𝟒 𝒏𝑨.
𝒎\textsuperscript{−𝟐}\textbf{\emph{. Todos los valores permanecen dentro
del régimen lineal dominado por conducción.}}
\end{quote}

El potencial en estado estacionario, la magnitud del campo eléctrico y
la magnitud de la densidad de corriente que rodean la corona inferior de
múltiples puntos se representan en la \textbf{Figura 4}. A la escala de
toda la corona (\textbf{Figura 4a}), cada uno de los cinco puntos
metálicos produce una depresión distinta en el campo de potencial. Estos
``embudos'' locales, típicamente 20-40 V más profundos que el aire
circundante a la misma altura, se fusionan en una perturbación más
amplia a unos pocos cm de la corona. El zoom cerca del punto central
(\textbf{Figura 4b}) muestra un canal de potencial estrecho y
verticalmente alargado cuyo mínimo se encuentra alrededor de 250-260 V,
correspondiendo a una caída de potencial lateral de 30-40 V sobre
escalas de mm-

a cm. Esto resulta tanto de la fuerte curvatura del vértice central como
del confinamiento lateral impuesto por los puntos vecinos de la corona.

La magnitud del campo eléctrico (\textbf{Figura 4c} y especialmente 4d)
exhibe cinco lóbulos de alto campo distintos, uno por vértice. Su
superposición forma un envoltorio estructurado en lugar de máximos
aislados, reflejando fuertes interacciones inter-puntos. Dentro del zoom
central (\textbf{Figura 4d}) el campo alcanza 𝐸\textsubscript{𝑚𝑎𝑥} ≈
80,9 𝑘𝑉. 𝑚\textsuperscript{−1}, casi 700 veces el campo ambiental de
tiempo despejado. Este campo cae a unos pocos kV/m dentro de
aproximadamente 1 cm, regresando a valores cercanos al ambiente dentro
de

varios cm. La amplificación permanece al menos un orden de magnitud por
debajo de los umbrales típicos de arranque de corona para conductores
afilados a presión de suelo, confirmando la operación en el régimen
pre-corona, óhmico.

Los mapas de densidad de corriente (\textbf{Figuras 4e} y \textbf{4f})
indican la estructura del campo eléctrico. Cada punto de la corona
genera un estrecho chorro de corriente que se difunde hacia abajo y se
fusiona con la corriente de conducción de fondo. El vértice central
nuevamente produce el valor más alto, 𝐽\textsubscript{max} ≈ 1,34 𝑛𝐴.
𝑚\textsuperscript{−2} que corresponde a un flujo microscópico similar al
estimado en el colector de varilla de punto único, \emph{es decir} 𝑁̇ ≈ 8
× 10\textsuperscript{9}𝑐ℎ𝑎𝑟𝑔𝑒𝑠. 𝑠\textsuperscript{−1}.
𝑚\textsuperscript{−2}. Dado que la conductividad varía solo débilmente
dentro de la región de 5-15 cm que rodea la corona, la estructura de
\textbar 𝐽\textbar{} está gobernada casi enteramente por el modelado
geométrico del campo en lugar de la estratificación atmosférica.

Bajo forzamiento de tiempo despejado, todos los campos simulados
permanecen bien por debajo de los umbrales de ruptura del aire,
confirmando la operación en un régimen lineal, óhmico pre-corona. La
amplificación del campo surge únicamente de la concentración geométrica
en los puntos metálicos, produciendo patrones suaves y monótonos de
campo eléctrico y densidad de corriente consistentes con 𝐽 = 𝜎𝐸. Estos
resultados definen una línea de base cuantitativa para potenciales,
factores de amplificación y chorros de conducción, contra los cuales se
puede comparar un forzamiento más fuerte. Ahora examinamos cómo responde
la geometría bajo condiciones idealizadas de tormenta, donde los campos
ambientales aumentan dramáticamente.

\subsection{Condiciones de tiempo
tormentoso}\label{condiciones-de-tiempo-tormentoso}

La respuesta de la varilla superior de punto único se estudia bajo un
forzamiento idealizado de tipo tormenta, usando los parámetros de la
Tabla 2. El campo eléctrico de fondo se aumenta a 𝐸\textsubscript{0,atm}
≈ 5000 𝑉. 𝑚\textsuperscript{−1} y la corriente de conducción impuesta
invierte dirección, alcanzando

𝐽\textsubscript{0,𝑎𝑡𝑚} ≈ +5 𝑛𝐴. 𝑚\textsuperscript{−2}. Esta
configuración proporciona un sustituto lineal para campos mejorados por
tormenta sin modelar explícitamente la carga de nubes, la carga espacial
o la estructura electrostática tridimensional de las nubes tormentosas.
Como consecuencia, los campos obtenidos cerca de las puntas representan
límites superiores sobre el estado pre-corona; una atmósfera real
transicionaría a una corona dominada por carga espacial antes de que se
mantengan tales campos.

La \textbf{Figura 5} muestra el potencial en estado estacionario, la
magnitud del campo eléctrico y la magnitud de la densidad de corriente
alrededor de la punta superior para este escenario de tipo tormenta. El
pozo de potencial alrededor del vértice (\textbf{Figura 5a}) es
significativamente más profundo y comprimido que bajo tiempo despejado,
reflejando tanto el forzamiento más fuerte como la corriente ascendente
más alta. El potencial local cerca del vértice de la punta única cae por
debajo de −12 𝑘𝑉, comparado con el rango de −260 𝑉 visto anteriormente,
indicando una intensificación de un orden de magnitud de los gradientes.
El campo eléctrico resultante (\textbf{Figura 5b}) alcanza un pico de
𝐸\textsubscript{𝑚𝑎𝑥} ≈ 735 𝑘𝑉. 𝑚\textsuperscript{−1} localizado en una
región sub-centimétrica. Este valor aún está por debajo de las
estimaciones canónicas de ruptura para aire a presión de suelo (∼ 2,5 −
2,8 𝑀𝑉. 𝑚\textsuperscript{−1}) pero cae dentro del rango inferior de
mediciones empíricas de arranque de corona para conductores afilados en
condiciones húmedas o pre-tormenta. Dentro de 1 cm del vértice, el campo
cae a unos pocos cientos de kV/m y regresa a valores de fondo dentro de
unos pocos cm, preservando el patrón de amplificación fuertemente
localizado observado en tiempo despejado.

La magnitud de la densidad de corriente correspondiente (\textbf{Figura
5c}) tiene un pico en

𝐽\textsubscript{𝑚𝑎𝑥} ≈ 5,4 µ𝐴. 𝑚\textsuperscript{−2} consistente con el
escalado 𝐽 = 𝜎𝐸, dada la conductividad elevada a la altitud de la punta.
Esto representa un aumento de casi tres órdenes de magnitud en relación
con los valores de tiempo despejado (∼ 1,36 𝑛𝐴. 𝑚\textsuperscript{−2}),
sin embargo la estructura espacial permanece como un estrecho chorro
axial que se difunde suavemente hacia la corriente ascendente de fondo.

A pesar de los grandes campos y corrientes absolutos, la solución
calculada retiene la firma cualitativa de un régimen de conducción
lineal: las superficies equipotenciales continúan convergiendo
monótonamente, \textbar 𝐸\textbar{} aumenta suavemente hacia el vértice
y no se forma bloqueo de campo ni meseta. Este comportamiento indica
que, dentro de las limitaciones de un modelo sin carga espacial, el
sistema permanece matemáticamente óhmico incluso en intensidades donde
una atmósfera real ya habría iniciado la corona. Los resultados deben,
por lo tanto, interpretarse como un estado pre-corona cercano al umbral,
con el inicio físico de la descarga inferido de criterios externos, no
de desviaciones dentro de los campos simulados.

La respuesta del distribuidor de corona de cinco puntos bajo forzamiento
de tipo tormenta se muestra en la \textbf{Figura 6}. Comparado con el
caso de tiempo despejado, el campo de fondo más fuerte y la corriente de
conducción invertida producen un pozo de potencial marcadamente más
profundo y confinado alrededor de cada vértice. A la escala de toda la
corona (\textbf{Figura 6a}), las superficies equipotenciales colapsan
firmemente alrededor de los cinco puntos metálicos, con potenciales
locales en la región de zoom (\textbf{Figura 6b}) cayendo alrededor de
−10 𝑘𝑉. Como en la varilla de punto único anterior, estos mínimos forman
canales estrechos y verticalmente alargados, pero el confinamiento
lateral impuesto por los picos vecinos afila los gradientes aún más.

La magnitud del campo eléctrico correspondiente (\textbf{Figuras 6c-6d})
muestra un conjunto de 5 lóbulos de alto campo. Su superposición produce
un envoltorio estructurado de campo amplificado, con el pico central
alcanzando 𝐸\textsubscript{max} ≈ 2845 𝑘𝑉. 𝑚\textsuperscript{−1}. Este
valor es aproximadamente 40 veces mayor que el máximo de la varilla de
punto único tipo tormenta y se encuentra directamente dentro del rango
empírico de ruptura para aire húmedo a nivel del suelo. También excede
las estimaciones típicas de arranque de corona difusa para conductores
afilados. Debido a que el modelo es estrictamente lineal y excluye
procesos de carga espacial o ionización, estos campos simulados deben
interpretarse como límites superiores del estado pre-descarga; una
atmósfera real comenzaría la formación de corona y el bloqueo de campo
antes de que se mantuvieran tales valores. Espacialmente, la
amplificación permanece confinada en una región sub-centimétrica
alrededor del vértice, decayendo al rango de cientos de kV/m dentro de
\textasciitilde1 cm y a niveles de fondo dentro de unos pocos cm.

Los mapas de densidad de corriente (\textbf{Figuras 6e-6f}) reproducen
esta estructura. Un chorro axial nitidamente enfocado emerge de cada
punto, con el central produciendo 𝐽\textsubscript{max} ≈ 5,4 µ𝐴.
𝑚\textsuperscript{−2}, más de tres órdenes de magnitud por encima de su
homólogo de tiempo despejado y consistente con el 𝐸\textsubscript{max}
elevado a través de 𝐽 = 𝜎𝐸. Aunque extremadamente grandes para un modelo
de solo conducción, el patrón espacial permanece suave y monótono, sin
el bloqueo de campo o las características no lineales que acompañarían
al inicio real de la corona.

En conjunto, estos resultados muestran que la corona de múltiples puntos
actúa como un concentrador de campo excepcionalmente eficiente bajo
forzamiento de tipo tormenta,

mucho más que el colector superior de punto único. Sin embargo, porque
el modelo descuida la carga espacial, la ionización y la física de
descarga, los picos predichos en escala MV/m representan límites
superiores idealizados en lugar de valores operativos realistas; en

la práctica, la corona, la actividad de streamers y la regulación de
carga asociada intervenirían bien antes de que se pudieran alcanzar los
campos simulados.

\begin{quote}
\includegraphics{./images/image5.jpeg}
\end{quote}

\begin{quote}
\textbf{\emph{Figura 5. Respuesta del colector de varilla de punto único
bajo forzamiento idealizado de tipo tormenta (}}𝑬𝟎,𝒂𝒕𝒎 ≈ 𝟓 𝒌𝑽.
𝒎\textsuperscript{−𝟏}\textbf{\emph{,}} 𝑱𝟎,𝒂𝒕𝒎 ≈ 𝟓 𝒏𝑨.
𝒎\textsuperscript{−𝟐}\textbf{\emph{). (a) Potencial V mostrando profunda
compresión de isopotenciales cerca del vértice metálico flotante. (b)
Magnitud del campo eléctrico}} \textbar 𝑬\textbar{}
\textbf{\emph{alcanzando}} 𝑬𝒎𝒂𝒙 = 𝟕𝟑𝟓 𝒌𝑽.
𝒎\textsuperscript{−𝟏}\textbf{\emph{, acercándose a rangos empíricos de
arranque de corona. (c) Magnitud de la densidad de corriente}}
\textbar 𝑱\textbar{} \textbf{\emph{con pico en}} 𝑱𝒎𝒂𝒙 = 𝟓. 𝟒 µ𝑨.
𝒎\textsuperscript{−𝟐}\textbf{\emph{. Los valores representan límites
superiores pre-corona dentro del modelo lineal sin carga espacial.}}
\end{quote}

\includegraphics{./images/image6.jpeg}

\begin{quote}
\textbf{\emph{Figura 6. Respuesta de la corona inferior de cinco puntos
bajo forzamiento idealizado de tipo tormenta (}}𝑬𝟎,𝒂𝒕𝒎 ≈ 𝟓 𝒌𝑽.
𝒎\textsuperscript{−𝟏}\textbf{\emph{,}} 𝑱𝟎,𝒂𝒕𝒎 ≈ 𝟓 𝒏𝑨.
𝒎\textsuperscript{−𝟐}\textbf{\emph{) con ángulo de vértice}} θ
\textbf{\emph{= 4°. (a,b) Potencial V mostrando un pozo colapsado
profundamente y fuerte confinamiento de isopotenciales alrededor de cada
vértice. (c,d) Magnitud del campo eléctrico}} \textbar 𝑬\textbar{}
\textbf{\emph{con un máximo localizado}} 𝑬𝒎𝒂𝒙 = 𝟐. 𝟖𝟒 𝑴𝑽.
𝒎\textsuperscript{−𝟏}\textbf{\emph{, acercándose a rangos empíricos de
ruptura. (e,f) Magnitud de la densidad de corriente}}
\textbar 𝑱\textbar{} \textbf{\emph{con pico en}} 𝑱𝒎𝒂𝒙 = 𝟐. 𝟖 µ𝑨.
𝒎\textsuperscript{−𝟐}\textbf{\emph{. Estos valores representan campos
superiores, pre-descarga dentro del modelo lineal sin carga espacial.}}
\end{quote}

\subsection{Influencia del ángulo de
vértice}\label{influencia-del-uxe1ngulo-de-vuxe9rtice}

Para evaluar la sensibilidad de la respuesta del electrovegetómetro a
laafilado del colector superior de Bertholon, se realiz un barrido
paramétrico para diferentes valores del ángulo de vértice θ de la
varilla de punto único, manteniendo todos los demás parámetros
geométricos y atmosféricos sin cambios. El ángulo θ se define aquí como
el ángulo de apertura completo de la punta triangular en la
configuración 2D (\textbf{Figura 2b}). Se exploraron valores desde 2°
hasta 90° (cuña roma de ángulo recto). Para cada θ, se extrajo la
magnitud máxima del campo eléctrico (Emax) del aire adyacente a la punta
superior y a la corona inferior. Además, la densidad de corriente de
conducción máxima (Jmax) también se extrajo en esas mismas regiones. Los
resultados se resumen en la \textbf{Figura 7}.

La \textbf{Figura 7a} muestra que el campo pico cerca de la varilla de
punto único disminuye monotónicamente a medida que la punta se rompe.
Bajo forzamiento de tiempo despejado (triángulos rojos), Emax cae desde
aproximadamente 20 kV/m para vértices muy afilados hasta 13 kV/m a 90°.
La dependencia es suave y relativamente débil: en todo el rango de
0-90°, la variación permanece dentro de un factor ≲1,5. La misma
tendencia se mantiene bajo forzamiento de tipo tormenta (triángulos
azules), con Emax disminuyendo desde 800 kV/m para puntas agudas hasta
450 kV/m para la cuña más roma. Por lo tanto, afilar

el vértice amplifica el campo local, pero solo al nivel de decenas de
porcentaje para ángulos realistas. Esta suave sensibilidad angular se
espera en una configuración de conductor flotante dominada por el
forzamiento de columna a gran escala: cambiar θ modifica la curvatura
local, por lo tanto el factor de amplificación geométrica, pero el
potencial general del metal se ajusta de manera autoconsistente para
satisfacer la corriente de conducción atmosférica impuesta. En otras
palabras, la filo de la punta controla cómo una caída de potencial casi
fija se concentra en gradientes de mm-cm, en lugar de cambiar el voltaje
global disponible para la estructura.

Para comparación, la \textbf{Figura 7a} también informa el campo máximo
en la corona de cinco puntos (estrellas). En ambos regímenes
atmosféricos, el campo pico de la corona permanece esencialmente
constante a medida que θ varía. En tiempo despejado, se mantiene cerca
de 80 kV/m y bajo forzamiento de tipo tormenta cerca de 2700-3000 kV/m.
Esta insensibilidad indica que la corona inferior está controlada casi
enteramente por su propia geometría de vértice y por el forzamiento
ambiental, con solo un acoplamiento despreciable a la filo del colector
superior distante. Prácticamente, esto significa que las incertidumbres
en la fabricación histórica de la punta superior no se propagan a la
amplificación de campo cercana al dosel producida por la corona
terminal.

\begin{quote}
\includegraphics{./images/image7.png}
\end{quote}

\begin{quote}
\textbf{\emph{Figura 7. (a) Campo eléctrico máximo}} 𝑬𝒎𝒂𝒙
\textbf{\emph{y (b) densidad de corriente}} 𝑱𝒎𝒂𝒙 \textbf{\emph{como
funciones del ángulo de vértice para el colector superior, bajo
forzamiento de tiempo despejado (rojo) y de tipo tormenta (azul). Los
valores de tiempo despejado muestran débil dependencia de la geometría
de la punta, mientras que las condiciones de tipo tormenta producen
fuerte amplificación geométrica. La corona inferior alcanza los campos
más altos (}}𝟐. 𝟖 𝑴𝑽. 𝒎\textsuperscript{−𝟏}\textbf{\emph{) y corrientes
(}}𝟓 𝒎𝑨. 𝒎\textsuperscript{−𝟐}\textbf{\emph{), representando
estimaciones superiores pre-descarga en el modelo óhmico lineal.}}
\end{quote}

Los máximos correspondientes de densidad de corriente de conducción se
muestran en la \textbf{Figura 7b}. Como predice la relación óhmica, Jmax
cerca de la punta superior sigue la misma disminución monótona con θ que
Emax. En tiempo despejado (triángulos rojos), la densidad de corriente
pico disminuye suavemente (nuevamente por menos de un factor de dos)
entre puntas afiladas y romas. En condiciones de tipo tormenta
(triángulos azules) ocurre la misma reducción, pero a valores que son
tres órdenes de magnitud más altos

que en tiempo despejado, reflejando el aumento impuesto tanto en el
campo de fondo como en la conductividad. En cuanto al campo eléctrico,
los máximos de corriente de la corona (estrellas) son esencialmente
planos con θ, confirmando que el comportamiento del distribuidor
inferior está desacoplado de la filo del colector en este régimen
pre-corona.

Una característica adicional de la \textbf{Figura 7b} es que, bajo
forzamiento de tipo tormenta, la punta superior exhibe un Jmax más
grande que la corona aunque la corona alcanza un Emax más alto. Esta
reversión surge de la estratificación de conductividad: el colector
superior opera a mayor altitud donde σ(z) es mayor, por lo que un campo
dado produce una corriente de conducción proporcionalmente más grande.
En contraste, el campo más fuerte de la corona se encuentra en una capa
de menor σ y, por lo tanto, lleva una densidad de corriente óhmica más
pequeña a pesar de su mayor amplificación geométrica.

En conjunto, el barrido del ángulo de vértice muestra que las
predicciones pre-corona para el colector superior son robustas a
incertidumbres plausibles en la filo de la punta. Si la punta metálica
del siglo XVIII era afilada como una aguja o moderadamente roma, cambia
Emax y Jmax solo por un factor modesto, no por órdenes de magnitud. Por
lo tanto, las conclusiones cualitativas extraídas de la geometría base
de θ

≈ 4° (a saber, que la operación de tiempo despejado permanece
seguramente sub-coronal mientras que el forzamiento de tiempo de
tormenta puede empujar la punta hacia rangos de arranque) no dependen
críticamente de una reconstrucción exacta de la forma de la punta. La
fuerte amplificación producida por la corona de múltiples puntos, por
otro lado, está gobernada por la propia filo y espaciado de la corona y
es efectivamente independiente del ángulo del colector. Finalmente, debe
recordarse que estas tendencias angulares se obtienen en una
aproximación de cuña 2D. Una punta cónica completamente 3D produciría
diferentes factores de amplificación absolutos, pero se espera que la
sensibilidad monótona y débil a θ persista, ya que refleja el ajuste de
potencial flotante y la localización de la caída de potencial en lugar
de cualquier particularidad de la geometría 2D.

\subsection{Influencia de la geometría del
conductor}\label{influencia-de-la-geometruxeda-del-conductor}

La Configuración I del electrovegetómetro parece efectiva bajo
condiciones de tiempo de tormento pero esencialmente inactiva en tiempo
despejado. Para explorar si geometrías alternativas podrían mejorar su
respuesta de tiempo despejado, se examinaron cuatro variantes del diseño
del conductor:

\begin{itemize}
\item
  La configuración I con colector de varilla de punto único y
  distribuidor de corona
\item
  La configuración II: el punto superior único se reemplaza por una
  pequeña corona de puntas afiladas montada en la cabeza del mástil,
  mientras que el distribuidor inferior permanece sin cambios.
\item
  La configuración I extendida: retiene un punto superior único pero
  extiende el mástil metálico aproximadamente 1 m, de modo que la punta
  ahora sobresale más alto en la columna atmosférica. La corona inferior
  y los soportes permanecen sin cambios. Esta
\end{itemize}

\begin{quote}
configuración no corresponde a un dibujo histórico específico pero sirve
para probar la sensibilidad del dispositivo a la altura del colector y a
los cambios asociados en el potencial y la conductividad de fondo.
\end{quote}

\begin{itemize}
\tightlist
\item
  La configuración de corona aislada: se retira todo el ensamblaje
  metálico de mástil-brazo para dejar solo la corona de cinco puntos a
  la altura del dosel. Esta configuración ``mínima'' proporciona un caso
  de control en el que la corona interactúa directamente con el campo
  ambiental sin ninguna guía o concentración por un conductor elevado.
\end{itemize}

Para las cuatro configuraciones, los valores máximos del campo eléctrico
obtenidos en la corona inferior se presentan en la \textbf{Figura 8a}
bajo forzamiento de tiempo despejado y de tipo tormenta. En tiempo
despejado, todas las configuraciones producen campos pico casi idénticos
del orden de 10\textsuperscript{2} kV.m\textsuperscript{−1}, indicando
que la geometría local de la corona domina la amplificación del campo y
que la presencia, forma o altura del colector superior tiene solo una
influencia menor en el régimen de campo débil. Bajo forzamiento de tipo
tormenta, sin embargo, las tres configuraciones con mástil (I, II e I
extendida) alcanzan todos campos pico comparables de aproximadamente
2,5-3 × 10\textsuperscript{3} kV.m\textsuperscript{−1}, mientras que la
configuración de solo corona permanece significativamente más baja
(aproximadamente la mitad de ese valor). Este contraste muestra que el
mástil metálico elevado canaliza eficientemente el potencial de columna
mejorado por tormenta hacia la región de la corona, pero que reemplazar
un colector de punto único por un colector de múltiples puntos o
extender modestamente su altura no cambia materialmente el campo máximo
local en el distribuidor.

Los máximos correspondientes de densidad de corriente de conducción,
𝐽\textsubscript{max}, se muestran en la \textbf{Figura 8c} en una escala
logarítmica. En condiciones de tiempo despejado, las cuatro
configuraciones se agrupan alrededor de 𝐽\textsubscript{max} ≈
10\textsuperscript{−3} 𝜇A.m\textsuperscript{−2}, confirmando que las
corrientes débiles y sub-coronas están esencialmente determinadas por la
corriente de conducción atmosférica ambiental y por la geometría de la
corona. Cuando se aplica forzamiento de tipo tormenta,

𝐽\textsubscript{max} aumenta casi tres órdenes de magnitud, alcanzando
unos pocos

𝜇A.m\textsuperscript{−2} para las tres configuraciones con mástil,
mientras que el caso de solo corona nuevamente permanece más bajo en un
factor de aproximadamente dos. El comportamiento paralelo de
𝐸\textsubscript{max} y 𝐽\textsubscript{max} en las cuatro
configuraciones es consistente con el escalado óhmico 𝐽 = 𝜎𝐸 y confirma
que el rol principal del mástil y el colector superior es aumentar la
cantidad de corriente de fondo interceptada y canalizada a través de la
corona, en lugar de introducir diferencias cualitativas fuertes entre
los diseños específicos del colector.

\begin{quote}
\includegraphics{./images/image8.jpeg}
\end{quote}

\begin{quote}
\textbf{\emph{Figura 8. (a) Cuatro geometrías distintas del
electrovegetómetro utilizadas en las simulaciones. Las regiones marrones
denotan soportes de madera, las líneas rojas conductores metálicos. (b)
Campo eléctrico máximo en el distribuidor de corona inferior, Emax, bajo
forzamiento de tiempo despejado y de tipo tormenta, mostrando que todas
las configuraciones con mástil producen campos pico similares mientras
que el caso de solo corona permanece más bajo en el régimen de tormenta.
(c) Densidad de corriente de conducción máxima correspondiente, Jmax, en
escala logarítmica, ilustrando el aumento de tres órdenes de magnitud de
condiciones despejadas a tormentosas y las corrientes mejoradas
obtenidas cuando está presente un mástil metálico elevado.}}
\end{quote}

\section{Discusión}\label{discusiuxf3n}

Las simulaciones presentadas anteriormente proporcionan un primer puente
cuantitativo entre la descripción cualitativa de Bertholon de su
electrovegetómetro y los conceptos modernos del circuito eléctrico
atmosférico global y la amplificación de campo pre-corona. Varios puntos
emergen con respecto a (i) la plausibilidad física de la narrativa
original de Bertholon, (ii) la magnitud esperada y localización de los
efectos eléctricos cerca del dosel y (iii) las limitaciones del marco de
modelado actual y sus implicaciones tanto para la interpretación
histórica como para las reclamaciones modernas de ``electrocultura''.

\subsection{Operación de tiempo despejado: amplificación localizada pero
modesta}\label{operaciuxf3n-de-tiempo-despejado-amplificaciuxf3n-localizada-pero-modesta}

Bajo forzamiento de tiempo despejado, el electrovegetómetro reconstruido
se comporta como una perturbación débil de la corriente de conducción
global. El colector superior y la corona inferior generan ambas fuertes
amplificaciones locales del campo (de dos a tres órdenes de magnitud en
relación con el campo ambiental de ≈10² V.m⁻¹) pero estas regiones
amplificadas permanecen confinadas en vecindarios de mm-cm alrededor de
las puntas. El punto superior único concentra la caída de potencial de
fondo en un embudo estrecho y alcanza campos pico del orden de unos
pocos ×10⁴ V.m⁻¹, mientras que la corona de cinco puntos alcanza
aproximadamente 8×10⁴ V.m⁻¹ cerca del vértice central.

Desde un punto de vista puramente electrodinámico, estos campos son lo
suficientemente altos como para distorsionar significativamente las
trayectorias iónicas y aumentar localmente las tasas de ionización y
adjunción, si la atmósfera ya está débilmente ionizada. Sin embargo, las
densidades de corriente óhmicas asociadas permanecen en el rango de
pA.m⁻² a nA.m⁻² bajo, solo

moderadamente por encima de la corriente de conducción de fondo. La
corriente total que puede canalizarse a través de una huella realista de
corona, por lo tanto, permanece extremadamente pequeña.

Para el microentorno de la planta, esto implica que, en clima sin
perturbaciones, el dispositivo podría plausiblemente crear estrechos
``puntos calientes'' de campo amplificado y flujo iónico inmediatamente
alrededor de cada punta, pero estas perturbaciones decaerían hasta
valores cercanos al fondo sobre distancias comparables o menores que las
escalas típicas de longitud de hoja y dosel. Cualquier forzamiento
electrodinámico directo sobre el tejido vegetal (a través de cargas
superficiales inducidas, flujos iónicos impulsados por campo en la capa
límite de la hoja o modificación de patrones de deposición de aerosoles)
sería, por lo tanto, altamente localizado y probablemente intermitente.

En consecuencia, el régimen de tiempo despejado parece compatible con la
idea cualitativa de Bertholon de una influencia continua pero ``suave''
de la electricidad atmosférica, sin embargo los resultados actuales
sugieren que esta influencia es físicamente sutil: el electrovegetómetro
no actúa como un fuerte inyector de carga o como un sumidero/fuente
significativo en el circuito global. Más bien, remodela la corriente de
conducción preexistente local en estrechos chorros sin aumentar
significativamente su magnitud.

\subsection{Forzamiento de tormenta y la plausibilidad de ``aigrettes''
luminosas}\label{forzamiento-de-tormenta-y-la-plausibilidad-de-aigrettes-luminosas}

Bajo forzamiento idealizado de tipo tormenta, la misma geometría opera
en un régimen muy diferente. Cuando el campo de fondo y la conductividad
se aumentan a valores representativos de condiciones pre-tormenta o
afectadas por tormenta, el punto superior y la corona inferior pueden
alcanzar campos eléctricos pico que van desde varios 10⁵ V.m⁻¹ hasta la
escala de MV.m⁻¹ en el modelo lineal. Estos niveles son comparables a o
exceden los umbrales empíricos de arranque de corona para conductores
afilados

a presión de suelo y se acercan a las estimaciones convencionales de
ruptura.

Es importante enfatizar que, porque el modelo descuida la carga espacial
y la retroalimentación de ionización, estos valores en escala MV.m⁻¹
deben interpretarse como \emph{límites superiores} del estado pre-corona
en lugar de como campos sostenidos realistas: en una atmósfera real, la
corona se activaría antes y la carga espacial resultante filtraría
parcialmente las puntas, limitando el crecimiento adicional del campo.
No obstante, el hecho de que la solución óhmica pre-corona se acerque
naturalmente a los umbrales de arranque de corona bajo forzamiento GEC
mejorado es significativo. Proporciona una justificación cuantitativa
para los informes de Bertholon de brillos luminosos tenues (``aigrettes
lumineuses'') en las puntas bajo clima perturbado y para las
observaciones históricas del fuego de San Elmo en mástiles de barcos y
agujas de iglesias.

En esta imagen, el electrovegetómetro puede actuar como un facilitador
de descargas de corona durante períodos cuando la columna atmosférica ya
está altamente estresada por la carga de tormenta cercana. El mástil
elevado asegura que el colector superior muestree una región de
ligeramente mayor conductividad y potencial, mientras que la corona
inferior de múltiples puntos concentra esta caída de potencial a la
altura del dosel. Las simulaciones muestran que las corrientes de tiempo
de tormenta pueden aumentar aproximadamente tres órdenes de magnitud en
comparación con el tiempo despejado mientras retienen perfiles de
conducción suaves y monótonos alrededor de las puntas. Una vez que se
alcanza el arranque de corona, la producción de iones y el transporte de
carga cerca del dosel podrían ser sustancialmente mayores que en el
régimen de tiempo despejado, potencialmente produciendo brillos visibles
y flujos iónicos mejorados.

Sin embargo, esta operación mejorada permanecería episódica, ligada a la
presencia de nubes electrificadas cercanas y puede no corresponder a la
``mejora'' benigna y en estado estacionario de la vegetación que
Bertholon tenía en mente. Además, la corona puede producir estrés
químico y mecánico localizado (ozono, NOx, UV, calentamiento a
micro-escala y cizalla de viento iónico) cuyos impactos en las plantas
son complejos y no necesariamente beneficiosos. El presente modelo no
puede abordar estos aspectos y, por lo tanto, no puede apoyar o refutar
directamente ningún beneficio agronómico reclamado de la operación de
tiempo de tormenta.

\subsection{Robustez a incertidumbres geométricas y el papel del
mástil}\label{robustez-a-incertidumbres-geomuxe9tricas-y-el-papel-del-muxe1stil}

Los barridos paramétricos indican que varios aspectos de la respuesta
del dispositivo son relativamente robustos a las incertidumbres
históricas en la geometría. Variar el ángulo de vértice del colector
superior en un amplio rango solo altera el campo máximo y la densidad de
corriente cerca de la punta por factores del orden de la unidad, no por
órdenes de magnitud. Esto refleja el hecho de que el potencial global
disponible para la estructura flotante está establecido por la columna
atmosférica impuesta, mientras que la filo de la punta controla cómo esa
caída de potencial se localiza en escalas de mm. Del mismo modo,
reemplazar el punto superior único por una pequeña corona o extender
modestamente el mástil, apenas afecta los campos pico en la corona
inferior en condiciones de tiempo despejado o de tormenta. Las
simulaciones muestran que, en el régimen de tormenta, todas las
configuraciones con mástil (Configuraciones I, II e I extendida)
producen Emax y Jmax similares en la corona a nivel del dosel, mientras
que una corona aislada sin mástil alcanza máximos significativamente más
bajos. Esto sugiere que el rol electrodinámico principal del mástil es
interceptar una

porción mejorada de la corriente y el potencial de la columna de fondo,
canalizándolos hacia el distribuidor, en lugar de codificar
sensiblemente los detalles del diseño del colector.

Desde una perspectiva de reconstrucción histórica, esto es
tranquilizador: incluso si el hardware real de Bertholon se desviaba de
la geometría estilizada utilizada aquí, las conclusiones cualitativas
sobre la operación de tiempo despejado versus de tormenta, la
localización de campos y la importancia relativa del colector versus el
distribuidor es improbable que cambien drásticamente. La corona de
múltiples puntos cerca del dosel de cultivos surge como la
característica dominante para la amplificación local del campo, mientras
que el mástil y el colector condicionan principalmente cuánto
``voltaje'' se entrega a ella.

\subsection{Limitaciones del modelo
actual}\label{limitaciones-del-modelo-actual}

A pesar de su utilidad como herramienta exploratoria, el presente modelo
2D, lineal, óhmico tiene varias limitaciones importantes que restringen
la interpretación de los resultados:

\begin{itemize}
\item
  El modelo asume un perfil de conductividad prescrito y la ley lineal
  de Ohm, sin tratamiento explícito de producción de carga, adjunción,
  recombinación o deriva más allá de lo que está codificado en σ(z). Una
  vez que los campos se acercan a los umbrales de arranque de corona,
  esta aproximación se rompe. Los campos pico simulados en el régimen de
  tipo tormenta deben, por lo tanto, verse como indicativos de dónde es
  probable el arranque, no como predicciones literales de valores
  sostenidos de MV.m⁻¹.
\item
  La configuración en sección transversal del electrovegetómetro
  representa una estructura infinitamente larga, lo que conduce a puntas
  de tipo cuña en lugar de puntas cónicas 3D realistas. Esto afecta los
  factores de amplificación locales y la topología de las líneas de
  campo cerca de los vértices. Una simulación 3D completa probablemente
  produciría diferentes valores absolutos de Emax y posiblemente
  regiones de campo alto más amplias, aunque las tendencias básicas
  (amplificación localizada, decaimiento decimétrico al fondo,
  insensibilidad relativa al ángulo de vértice del colector) deberían
  permanecer cualitativamente similares.
\item
  Los soportes de madera se modelan con conductividades fijas y bajas y
  el suelo se trata implícitamente a través de condiciones de contorno
  en lugar de como un conductor estratificado y dependiente de la
  humedad. En la práctica, sin embargo, Bertholon dedicó considerable
  atención a mejorar el aislamiento: la sección enterrada del mástil se
  secó al fuego, se alquitranó, se envolvió en polvo de carbón y
  cemento, luego se asentó en albañilería, mientras que la parte sobre
  el suelo se pintó o cubrió con betún; la cabeza del mástil se protegió
  aún más con inserciones saturadas en resina y mangas de
  vidrio/mastique y los brazos horizontales se llevaron sobre caballetes
  aislantes con cuerdas de seda (Tabla 1). Estas medidas probablemente
  redujeron las fugas y aumentaron la resistencia del soporte, sin
  embargo sus propiedades aún dependían de la humedad, el envejecimiento
  y la contaminación, por lo que las rutas de fuga y el potencial
  flotante (y por lo tanto Emax y Jmax) pueden haber cambiado solo por
  factores moderados.
\item
  Los regímenes de ``tiempo despejado'' y ``tipo tormenta'' se
  representan perfiles simples y horizontales uniformes con J0,atm y
  E0,atm prescritos. Los entornos eléctricos atmosféricos reales sobre
  campos agrícolas están influenciados por la topografía, estructuras
  cercanas, carga de aerosoles e inhomogeneidades horizontales en la
  carga de tormenta. El campo local en el electrovegetómetro podría, por
  lo tanto, desviarse
\end{itemize}

\begin{quote}
significativamente de los perfiles idealizados adoptados aquí,
especialmente bajo condiciones convectivas complejas.
\end{quote}

\begin{itemize}
\tightlist
\item
  Las simulaciones no incluyen representaciones explícitas de
  vegetación. Las hojas, tallos y estructura del dosel pueden
  distorsionar el campo y actuar como conductores adicionales o
  aislantes parciales, redistribuyendo corrientes y modificando el
  paisaje de potencial cercano a la superficie. También pueden
  participar en el intercambio de carga a través de películas de
  conducción superficial, aberturas estomatales o capas de humectación.
  Los resultados actuales proporcionan, por lo tanto, una imagen de
  primer orden del campo ``en el aire'', no una descripción completa del
  acoplamiento planta-campo.
\end{itemize}

\subsection{Implicaciones para la historia y la práctica moderna de la
``electrocultura''}\label{implicaciones-para-la-historia-y-la-pruxe1ctica-moderna-de-la-electrocultura}

Dentro de estas limitaciones, los resultados del modelado apoyan una
reevaluación matizada del electrovegetómetro de Bertholon. Por un lado,
el dispositivo es físicamente capaz de generar fuertes amplificaciones
locales de campo y, bajo condiciones de tiempo de tormenta, de acercarse
o exceder los umbrales de arranque de corona, haciendo que las
``aigrettes'' luminosas sean plausibles. Por otro lado, en su estado
operativo principal de tiempo despejado (que dominaría la exposición
diaria), la estructura parece actuar principalmente como un concentrador
pasivo de campo con corrientes totales pequeñas y perturbaciones
altamente localizadas.

Esta imagen dual sugiere que la intuición conceptual de Bertholon sobre
la existencia de una influencia eléctrica persistente sobre la
vegetación no estaba completamente equivocada, pero que sus expectativas
sobre la fuerza y las consecuencias agronómicas de esta influencia
podrían haber sido optimistas. El electrovegetómetro es mejor
interpretado como una sonda delicada del circuito eléctrico global en la
capa cercana a la superficie y como un demostrador de concentración de
campo y fenómenos de corona, que como un ``mejorador'' agrícola robusto
en el sentido de ingeniería moderno.

Para las discusiones contemporáneas de ``electrocultura'', un término
que ha resurgido recientemente en contextos populares y a veces
pseudocientíficos, el presente trabajo subraya la necesidad de un
modelado cuidadoso y cuantitativo y experimentos controlados. Los
dispositivos que se asemejan a electrovegetómetros históricos pueden
efectivamente modificar el entorno eléctrico local alrededor de las
plantas, pero la magnitud, el signo y la relevancia biológica de estas
modificaciones no pueden inferirse solo de analogías cualitativas.
Requieren una combinación de mediciones atmosférico-eléctricas, estudios
detallados de fisiología vegetal y modelos que se extiendan más allá del
régimen pre-corona, lineal considerado aquí.

\section{Conclusión}\label{conclusiuxf3n}

Reexaminando el electrovegetómetro del siglo XVIII de Bertholon dentro
del marco de la electrodinámica atmosférica moderna, este estudio ha
desarrollado un modelo cuasi-estacionario, bidimensional, resistivo del
sistema aire-dispositivo, en el que la atmósfera se trata como un medio
óhmico que transporta la corriente de conducción global y la estructura
metálica se representa como un conductor flotante soportado por
aislantes de madera con fugas. La resolución de la ecuación de
conducción estacionaria con conductividad dependiente de la altitud

y configuraciones realistas ha producido estimaciones cuantitativas de
potenciales, campos eléctricos y densidades de corriente alrededor del
electrovegetómetro bajo forzamiento idealizado de tiempo despejado y de
tipo tormenta.

Los hallazgos principales pueden resumirse de la siguiente manera:

\begin{itemize}
\item
  Bajo condiciones típicas de tiempo despejado, las puntas afiladas del
  electrovegetómetro amplifican el campo eléctrico local entre dos y
  tres órdenes de magnitud en relación con el campo atmosférico de
  fondo, alcanzando decenas de kV.m⁻¹ en el punto superior único y
  aproximadamente 80 kV.m⁻¹ en la corona inferior de múltiples puntos.
  Las densidades de corriente asociadas permanecen en el rango de
  pA-nA/m2, consistente con una conducción débil en un medio ligeramente
  ionizado. Las perturbaciones de campo y corriente están fuertemente
  localizadas y decaen hasta valores cercanos al ambiente dentro de cm.
\item
  Cuando el campo de fondo y la conductividad se aumentan a valores
  representativos de condiciones pre-tormenta o afectadas por tormenta,
  la misma geometría produce campos pico en el rango de 10⁵-10⁶ V.m⁻¹.
  Estos niveles se acercan o exceden los umbrales empíricos de arranque
  de corona y se acercan a las estimaciones de ruptura, particularmente
  en la corona inferior. Dentro de los límites de una aproximación
  lineal sin carga espacial, estos resultados indican que el dispositivo
  de Bertholon podría haber producido brillos de corona y flujos iónicos
  mejorados durante clima perturbado, proporcionando una base física
  plausible para los informes históricos de ``aigrettes'' luminosas.
\item
  Las variaciones paramétricas del ángulo de vértice superior y la
  geometría del colector muestran que los campos pico y corrientes de
  tiempo despejado y de tormenta cerca de la corona inferior son en gran
  medida insensibles a los detalles finos del colector, siempre que esté
  presente un mástil elevado. La función principal del mástil es
  interceptar una porción del potencial de columna atmosférica y
  canalizarla hacia el distribuidor, mientras que la corona de múltiples
  puntos cerca de la altura del dosel permanece como la característica
  dominante que controla la amplificación local del campo.
\item
  Los resultados dependen de los supuestos de una geometría 2D,
  conducción óhmica lineal, perfiles de conductividad prescritos y
  propiedades de materiales simplificadas. Deben, por lo tanto, verse
  como límites superiores pre-corona en lugar de predicciones
  completamente realistas del comportamiento de descarga. Dentro de
  estos límites, el estudio sugiere que el electrovegetómetro de
  Bertholon era físicamente capaz de concentrar la electricidad
  atmosférica y ocasionalmente producir corona visible bajo forzamiento
  de tiempo de tormenta, mientras que su influencia de tiempo despejado
  sobre las plantas era probablemente sutil, localizada y difícil de
  cuantificar con instrumentos del siglo XVIII.
\end{itemize}

Finalmente, este trabajo ilustra cómo el modelado numérico moderno del
circuito eléctrico atmosférico global, incluso en una forma simplificada
pre-corona, puede afianzar nuestra comprensión de dispositivos
electrotécnicos históricos y sus modos plausibles de operación. Las
extensiones futuras podrían incluir simulaciones 3D completas,
tratamiento explícito de carga espacial y dinámica de corona,
representaciones más realistas de suelo y vegetación y experimentos
acoplados en entornos atmosférico-eléctricos controlados. Tales estudios
no solo refinarían la evaluación retrospectiva del electrovegetómetro de
Bertholon sino que también informarían los debates contemporáneos sobre
la significación física y biológica de las perturbaciones
atmosférico-eléctricas débiles en entornos agrícolas.

\section{Referencias}\label{referencias}

{[}Aplin, 2008{]} K. L. Aplin, R. G. Harrison, M. J. Rycroft.
Investigating Earth's Atmospheric Electricity: a Role Model for
Planetary Studies. Space Sci Rev, Vol. 137, pp.~11-27 (2008).
{[}https://doi.org/10.1007/s11214-008-9372-x{]}

{[}Winkler et al., 1898{]} J. H. Winkler, B. Franklin, T. F. Dalibard,
L. G. Le Monnier. Ueber Luftelektricität. Neudruck Von Schriften Und
Karten Über Meteorologie Und Erdmagnetismus Herausgegeben Von Professor
Dr.~G. Hellmann. N°11. Berlin, A. Asher \& Co.~(1898).

{[}Bailey, 2001{]} A. G. Bailey. The charging of insulator surfaces.
Journal of Electrostatics. Vol. 51-52, pp.~82-90 (2001).
{[}https://doi.org/10.1016/S0304-3886(01)00106-1{]}00106-1)

{[}Bertholon, 1783{]} Bertholon. De l'électricité des végétaux. A Paris,
chez P. F. Didot Jeune, quai des Augustins (1783).

{[}Guha, 2010{]} A. Guha, B. K. De, S. Gurubaran, S. S. De, K. Jeeva.
First results of fair-weather atmospheric electricity measurements in
Northeast India. J. Earth Syst. Sci., Vol. 119, pp.~221-228 (2010).
{[}https://doi.org/10.1007/s12040-010-0014-9{]}

{[}Herbert, 2012{]} K. B. H. Herbert. John Canton --- Pioneer
investigator of atmospheric electricity. Weather, RMetS, Vol. 52, Issue
9, pp.~286-290 (1997).
{[}https://doi.org/10.1002/j.1477-8696.1997.tb06326.x{]}

{[}Kudintseva, 2016{]} I. G. Kudintseva, A. P. Nickolaendro, M. J.
Rycroft, A. Odzimek. AC and DC global electric circuit properties and
the height profile of atmospheric conductivity. Annals of geophysics,
Vol. 59, 5, A0545 (2016). {[}https://www.doi.org/10.4401/ag-6870{]}

{[}Kundt, 1999{]} W. Kundt, G. Thuma. Geoelectricity: atmospheric
charging and thunderstorms. Journal of Atmospheric and Solar-Terrestrial
Physics, Vol. 61, pp.~955-963 (1999).
{[}https://doi.org/10.1016/S1364-6826(99)00061-9{]}00061-9)

{[}Mareev, 2014{]} E. A. Mareev, E. M. Volodin. Variation of the global
electric circuit and Ionospheric potential in a general circulation
model. Geophysical Research Letters, Vol. 41, Issue 24, pp.~9009-9016
(2014). {[}https://doi.org/10.1002/2014GL062352{]}

{[}McGinness, 2025{]} B. P. S. McGinness, R. Giles Harrison, K. L.
Aplin,

M. W. Airey. The role of point discharge in the historical development
of atmospheric electricity. Hist. Geo Space Sci., 16, 51-63 (2025).
{[}https://doi.org/10.5194/hgss-16-51-2025{]}

{[}Standler, 1979{]} R. B. Standler, W. P. Winn. Effects of coronae on
electric fields beneath thunderstorms. Quaterly J. of Royal
Meteorological Society, Vol. 105, 443, pp.~285-302 (1979).
{[}https://doi.org/10.1002/qj.49710544319Ci{]}

{[}Tinsley, 2007{]} B. A. Tinsley, G. B. Burns, L. Zhou. The role of the
global electric circuit in solar and internal forcing of clouds and
climate. Advances in Space Research Vol. 40, Issue 7, pp.~1126-1139
(2007). {[}https://doi.org/10.1016/j.asr.2007.01.071{]}

{[}Williams, 2009{]} E. R. Williams. C.T.R. Wilson versus G.C. Simpson:
Fifty years of controversy in atmospheric electricity. Atmospheric
Research; Vol. 91, Issues 2-4, pp.~259-271 (2009).
{[}https://doi.org/10.1016/j.atmosres.2008.03.024{]}

\end{document}
